% DER_XIIb-b: Grothendieck's text with the commentary set in, then the overview, inside a group that applies the chunk's own definitions
\dgrpart{1}{XIIb-b}{Section 5, part 2 of 2, manuscript pages 123--169}{XIIb-b\enspace\textperiodcentered\enspace Chapter XII, section 5, part 2 of 2\enspace\textperiodcentered\enspace manuscript pages 123--169}{Chapter XII \textperiodcentered\ \textup{XIIb-b}}
\dgrchunkbegin
% DER_XIIb-b: the definitions of the chunk's preamble, applied inside its group
\providecommand\IfDocumentMetadataT[1]{\IfDocumentMetadataTF{#1}{}}
\let\printedpage\relax\newcommand{\printedpage}[1]{\marginpar{\footnotesize p.~#1}}
\let\msp\relax\newcommand{\msp}[1]{\marginpar{\footnotesize [#1]}}
\let\Hot\relax\newcommand{\Hot}{\mathrm{Hot}}
\let\HOT\relax\newcommand{\HOT}{\mathrm{HOT}}
\let\Cat\relax\newcommand{\Cat}{\mathrm{Cat}}
\let\cX\relax\newcommand{\cX}{\mathcal{X}}
\let\sX\relax\newcommand{\sX}{\mathscr{X}}
\let\tX\relax\newcommand{\tX}{\tilde{\mathscr{X}}}
\let\Hom\relax\newcommand{\Hom}{\mathrm{Hom}}
\let\Fl\relax\newcommand{\Fl}{\mathrm{Fl}}
\let\tsX\relax\newcommand{\tsX}{\tilde{\mathscr{X}}}
\let\circled\relax\newcommand{\circled}[1]{\tikz[baseline=(char.base)]{\node[shape=circle,draw,inner sep=1pt] (char) {#1};}}
\let\uHom\relax\newcommand{\uHom}{\underline{\mathrm{Hom}}}
\let\tld\relax\newcommand{\tld}{^{\tilde{}}}
\let\hto\relax\newcommand{\hto}[1][]{\overset{#1}{\lhook\joinrel\longrightarrow}}
% Transcription repair: complete the truncated macro; the printed integral diagrams are unchanged.
\let\intdiag\relax\newcommand{\intdiag}[3]{\int\begin{tikzcd}[row sep=small,column sep=small,ampersand replacement=\&,baseline=(\tikzcdmatrixname-2-1.base)] \& #2 \\ #1 \arrow[ur]\arrow[dr] \& \\ \& #3 \end{tikzcd}}
\let\Ob\relax\newcommand{\Ob}{\operatorname{Ob}}
\let\Fibdisc\relax\newcommand{\Fibdisc}{\underline{\mathrm{Fib\ disc}}}
\let\cercle\relax\newcommand{\cercle}[1]{\tikz[baseline=(c.base)]\node[draw,circle,inner sep=1pt](c){#1};}

% ---- from Der-ChXIIb-012_p041-045.tex
\printedpage{83}
% head: A. Grothendieck : Dérivateurs, Ch. XII, Édition des Universités de Montpellier II et Paris VII
d'après le calcul connu du foncteur $\pi_!$, adjoint à gauche du foncteur $\pi^*$
\[
\HOT_W(e) \overset{\text{déf}}{=} \Hot_W \longrightarrow \HOT_W(\Phi)
\]
(\footnote[84]{Cela suppose qu'on a W(2), non seulement W(2 bis).}). Or pour tout objet $I$ dans Cat, de projection absolue $p_I : I \longrightarrow e$, le foncteur
\[
p_{I\,!} : \HOT_W(I) \longrightarrow \HOT_W(e) = \Hot_W
\]
doit être regardé comme l'équivalent `homotopique' du foncteur limite inductive $\varinjlim\limits_{I}$ des diagrammes de type $I^{\circ}$ dans la catégorie $\Hot_W$ - lesquelles limites inductives \textit{n'existent pas} dans $\Hot_W$ (sauf
\msp{page 123}
pour des catégories très particulières - les seules que je connaisse sont les catégories discrètes (qu'il faut peut-être même supposer finies en plus)).

Notons qu'on a un diagramme de foncteurs canoniques
\[
\begin{tikzcd}[column sep=tiny, row sep=small, ampersand replacement=\&]
 \& X_0 \arrow[dl, "\alpha_1"'] \arrow[dr, "\alpha_2"] \& \\
X_1 \arrow[dr, "i_1"'] \& {\shortstack{\scriptsize non\\ \scriptsize commuta-\\ \scriptsize tif!}} \& X_2 \arrow[dl, "i_2"] \\
 \& \cX, \&
\end{tikzcd}
\tag*{$(*)$}
\]
lequel diagramme n'a rien de commutatif dans Cat. \textit{Il est cependant commutatif dans} $\Hot_W$ (quel que soit le localiseur fondamental $W$), et même dans Cat\,hot, car si on désigne par
\[
i_0 : X_0 \longrightarrow \cX
\]
l'inclusion canonique, on a des homomorphismes canoniques
\[
i_1\alpha_1 \longrightarrow i_0 \longleftarrow i_2\alpha_2,
\]
ce qui montre que $i_1\alpha_1$ et $i_2\alpha_2$ sont homotopes, donc définissent la même flèche dans $\Hot_W$.

\msp{page 124}
Considérons dans $\Phi$ les deux fermés isomorphes à $\boldsymbol{\Delta}^1$
\[
\begin{aligned}
F_1 &= 1\backslash\Phi &&= \{0,1\}\\
F_2 &= 2\backslash\Phi &&= \{0,2\},
\end{aligned}
\]
on a
\[
F_1\cap F_2 = 0\backslash\Phi = \{0\}, \qquad F_1\cup F_2 = \Phi.
\]

\dgrblock{dgrb460}{A circle hidden by an ordinary pushout}
\begin{dcomment}{examples}{p123}{A circle hidden by an ordinary pushout}
The distinction between a homotopy colimit and a colimit in the homotopy category is already visible in a span of finite discrete categories.  Take $X_0=\{a,b\}$ and $X_1=X_2=e$, with both structure maps the unique functor.  Call the two side objects $l,r$.  Formula (A4) gives precisely four nonidentity arrows:
\[
 l\to a,\quad l\to b,\quad r\to a,\quad r\to b.
\]
No two nonidentity arrows are composable.  The nondegenerate nerve is consequently a graph with four vertices and four edges, namely a four-cycle.  Its realization is a circle.  One can compute this without any homotopy-colimit theorem: choose three of the edges as a spanning tree and collapse that tree; the remaining edge is one loop.  Thus $\pi_1(B\mathcal X)\cong\mathbb Z$, $H_1(B\mathcal X;\mathbb Z)\cong\mathbb Z$, and all reduced homology in other degrees vanishes.

The ordinary pushout $e\amalg_{\{a,b\}}e$ in $\Cat$ is $e$.  Indeed, a strictly compatible pair of functors from its two copies of $e$ to $Y$ chooses a single object of $Y$; both points of $X_0$ impose the same equality.  The canonical comparison $\mathcal X\to e$ therefore collapses the circle and is not in $W_\infty$.  This example will fail Definition~3 exactly where it should.

The ordinary pushout of the same diagram also exists in the homotopy category of spaces and is a point.  For a space $T$, a homotopy class $e\to T$ is a component of $T$.  Compatibility of two such classes along the discrete two-point space forces them to be the same component, so the set of compatible pairs is $\pi_0(T)$, represented by $e$.  A homotopy pushout instead remembers two paths implementing the two compatibilities; their difference produces the circle.  Thus even existence of the ordinary colimit does not give the desired answer.

The parenthetical uncertainty on manuscript p.~123 about infinite discrete indexing categories should not be carried into the modern statement: the homotopy category of spaces has all small coproducts, represented by disjoint unions of CW models.  What fails is a general ordinary-colimit calculus that recovers coherent homotopy colimits.  The four-object calculation isolates that failure without making an unnecessary nonexistence assertion.
\end{dcomment}


\printedpage{84}
% head: A. Grothendieck : Dérivateurs, Ch. XII, Édition des Universités de Montpellier II et Paris VII
Soient
\[
\begin{aligned}
1\backslash\cX &= \cX\times_{\Phi} 1\backslash\Phi &&= \cX\times_{\Phi} F_1\\
2\backslash\cX &= \cX\times_{\Phi} 2\backslash\Phi &&= \cX\times_{\Phi} F_2
\end{aligned}
\]
les images inverses de $F_1$, $F_2$ dans $\cX$, on a donc
\[
1\backslash\cX\cap 2\backslash\cX = X_0, \qquad 1\backslash\cX\cup 2\backslash\cX = \cX,
\]
donc le diagramme
\[
\begin{tikzcd}[column sep=small, row sep=small, ampersand replacement=\&]
 \& X_0 \arrow[dl, hook', "\alpha'_1"'] \arrow[dr, hook, "\alpha'_2"] \& \\
1\backslash\cX \arrow[dr, hook, "i'_1"'] \& \& 2\backslash\cX \arrow[dl, hook', "i'_2"] \\
 \& \cX, \&
\end{tikzcd}
\tag*{$(**)$}
\]
déduit du diagramme similaire
\[
\begin{tikzcd}[column sep=small, row sep=small, ampersand replacement=\&]
 \& \{0\} \arrow[dl, hook'] \arrow[dr, hook] \& \\
F_1 \arrow[dr, hook] \& \& F_2 \arrow[dl, hook'] \\
 \& \Phi, \&
\end{tikzcd}
\]
est (comme ce dernier) un carré d'immersions fermées à la fois cartésien et cocartésien. D'autre part, on a un homomorphisme
\msp{page 125}
canonique du diagramme $(*)$ dans $(**)$, qui est l'identité sur les sommets $X_0$ et $\cX$, et qui sur $X_1$, $X_2$ est donné par les inclusions canoniques
\[
\underbrace{X_1}_{\text{\scriptsize $\cX_1$ fibre de $\cX$ en 1}} \longrightarrow 1\backslash\cX, \qquad
\underbrace{X_2}_{\text{\scriptsize $\cX_2$ fibre de $\cX$ en 2}} \longrightarrow 2\backslash\cX.
\]
Ces inclusions, on le sait, sont $W$-asphériques (elles ont même un adjoint à droite), a fortiori elles sont dans $W$. Donc le diagramme non commutatif $(*)$ est isomorphe, en tant que diagramme dans $\Hot_W$, au diagramme commutatif cartésien et cocartésien $(**)$ formé d'immersions fermées. Comme $\cX$ est bel et bien somme amalgamée de ses deux fermées % [sic] fermées -> fermés, cf. « les deux fermés » p. 83
$1\backslash\cX$, $2\backslash\cX$ sous leur intersection $X_0$, cela étaye l'intuition qui voit $\cX$ comme une `somme amalgamée homotopique'. Donnons tout de suite l'application suivante :

\dgrblock{dgrb461}{Closed pieces give a genuine gluing diagram}
\begin{dcomment}{exposition}{p125:cela étaye l'intuition qui voit}{Closed pieces give a genuine gluing diagram}
Grothendieck now replaces the informal square by a strict cover, so that an excision axiom can actually be applied.  The construction is a categorical version of replacing the two legs of a span by mapping cylinders.

Let $C_k$ be the full subcategory of $\mathcal X$ on $X_k\amalg X_0$.  This is $k\backslash\mathcal X$ in the manuscript.  It is a cosieve: an arrow starting in its side fibre can only remain there or enter $X_0$, and arrows starting in $X_0$ remain in $X_0$.  We have
\[
 C_1\cap C_2=X_0,\qquad C_1\cup C_2=\mathcal X,
 \qquad \mathcal X=C_1\amalg_{X_0}C_2.
 \tag{A5}
\]
The last equality holds in $\Cat$, not just on objects.  A string of composable arrows cannot pass from one side fibre to the other; it can enter the common fibre but cannot leave it.  Giving a functor out of $\mathcal X$ is therefore exactly giving functors out of the two pieces that agree on $X_0$.  The intersection statement also makes the square cartesian.

The inclusion $j_k:X_k\to C_k$ has right adjoint $r_k$, defined by $r_k(a)=a$ on $X_k$ and $r_k(x)=\alpha_kx$ on $X_0$.  On the mixed morphisms it uses (A4).  Indeed, for every $a\in X_k$ and $z\in C_k$,
\[
 \operatorname{Hom}_{C_k}(j_ka,z)
       =\operatorname{Hom}_{X_k}(a,r_kz).
\]
The unit is the identity and the counit is the canonical transition arrow on the middle fibre.  Thus $Bj_k$ and $Br_k$ are mutually inverse homotopy equivalences, and both functors belong to every fundamental localizer.

For spaces, (A5) proves the promised homotopy-pushout calculation.  Every simplex of $N\mathcal X$ lies in $NC_1$ or $NC_2$, and their intersection is $NX_0$.  Hence it is the pushout of these nerves.  The two inclusions of $NX_0$ are simplicial monomorphisms, whose realizations are CW-subcomplex inclusions, so this ordinary union computes the homotopy pushout of $BC_1\leftarrow BX_0\to BC_2$.  The retractions $r_k$ give a strictly commuting map of this span to $BX_1\leftarrow BX_0\to BX_2$, because $r_k|_{X_0}=\alpha_k$; all three components are homotopy equivalences.  Homotopy invariance of the double mapping cylinder proves the result.

The phrase ``homomorphism of diagrams'' on p.~125 must retain this distinction: the inclusions from the original side fibres form a comparison with specified natural transformations, and become a comparison of commutative squares in $\Hot_W$.  The retractions just written give a strict map of spans in the other direction.  No strict commutativity of the original square is being asserted or needed.  For a general $W$, the same replacement reduces Proposition~13 to precisely the closed-cover form of the stated axiom W(7 bis).
\end{dcomment}


\textbf{Proposition 13.} \textit{Soit un diagramme} $\begin{tikzcd}[column sep=small, row sep=tiny, ampersand replacement=\&, baseline=(\tikzcdmatrixname-2-1.base)]
 \& X_1 \\
X_0 \arrow[ur, "\alpha_1"] \arrow[dr, "\alpha_2"'] \& \\
 \& X_2
\end{tikzcd}$ \textit{dans} Cat\textit{, d'où le diagramme} $(*)$\textit{, p. 123,}

\printedpage{85}
% head: A. Grothendieck : Dérivateurs, Ch. XII, Édition des Universités de Montpellier II et Paris VII
\[
\begin{tikzcd}[column sep=small, row sep=small, ampersand replacement=\&]
 \& X_0 \arrow[dl, "\alpha_1"'] \arrow[dr, "\alpha_2"] \& \\
X_1 \arrow[dr, "i_1"'] \& \& X_2 \arrow[dl, "i_2"] \\
 \& \cX, \&
\end{tikzcd}
\]
\msp{page 126}
\textit{commutatif dans} $\Hot_W$ \textit{(mais non dans} Cat\textit{, sauf si $X_0 = \emptyset$). Supposons que $W$ satisfait} W(7 bis) (axiome du carré cocartésien, cf. p. 75). \textit{Alors on a}
\[
\left\{
\begin{aligned}
\alpha_1\in W &\iff i_2\in W\\
\alpha_2\in W &\iff i_1\in W. \text{ (\footnotemark[85])}
\end{aligned}
\right.
\]
\footnotetext[85]{\textbf{NB} Pour l'implication $\Longrightarrow$ on n'a besoin que de la forme directe de W(7 bis), i.e. (avec les notations de XVI) de Loc(6$'$a).}
\textsc{Démonstration}. Utilisant l'homomorphisme précédent de $(*)$ dans $(**)$, on voit que chacune des flèches $\alpha_1$, $\alpha_2$, $i_1$, $i_2$ est dans $W$, si et seulement si la flèche correspondante $\alpha'_1$, $\alpha'_2$, $i'_1$ ou $i'_2$ l'est. Donc les équivalences ci-dessus équivalent à
\[
\begin{aligned}
\alpha'_1\in W &\iff i'_2\in W\\
\alpha'_2\in W &\iff i'_1\in W,
\end{aligned}
\]
lesquelles sont vraies en vertu de W(7 bis) (formulation duale) appliqué à $\cX$ et aux deux fermés $1\backslash\cX$, $2\backslash\cX$ de $\cX$.

\dgrblock{dgrb462}{Why converse excision is an additional axiom}
\begin{dcomment}{development}{p126:lesquelles sont vraies en vertu de W(7 bis)}{Why converse excision is an additional axiom}
Proposition~13 is correctly proved under its stated hypothesis.  The strength of that hypothesis nevertheless matters: the displayed equivalences contain converse excision, whereas the usual invariance of a homotopy pushout under weak equivalences gives only the forward implication.  A suspension makes the difference precise.

\textbf{Proposition (the acyclic test for converse excision).}\ours
Let $W$ contain the weak homotopy equivalences of small categories and satisfy two-out-of-three.  Suppose that, whenever $C=C_1\cup C_2$ is a union of two full cosieves with intersection $A$, the implication
\[
 C_2\longrightarrow C\text{ belongs to }W
 \quad\Longrightarrow\quad A\longrightarrow C_1\text{ belongs to }W
 \tag{A6}
\]
holds.  Then every category $A$ whose classifying space is connected and integrally acyclic is $W$-aspherical.

\emph{Proof.}
Adjoin to $A$ two distinct initial vertices $l,r$, with one arrow from each vertex to every object of $A$, and no arrows between $l$ and $r$.  Call the resulting category $C$ and its two one-vertex enlargements $C_1,C_2$.  The full subcategories $C_1,C_2$ are cosieves, and their intersection is $A$.  Their classifying spaces are cones on $BA$; the classifying space of $C$ is their union along $BA$, namely the unreduced suspension $\Sigma BA$.

Since $BA$ is connected, van Kampen applied to the two slightly enlarged cones makes $\Sigma BA$ simply connected.  The reduced Mayer--Vietoris sequence gives $\widetilde H_n(\Sigma BA;\mathbb Z)\cong\widetilde H_{n-1}(BA;\mathbb Z)=0$ for every $n$.  A simply connected acyclic CW complex is contractible: if it had a first nonzero homotopy group, the Hurewicz theorem would identify it with a nonzero homology group; with all homotopy groups zero, Whitehead's theorem applies.  Equivalently, use \citet[Corollary~4.33]{Hatcher2002}.  Thus $C_2\to C$ is a weak homotopy equivalence, and lies in $W$.  By (A6), $A\to C_1$ lies in $W$.  Since $C_1\to e$ also lies in $W$, two-out-of-three gives $A\to e\in W$.\hfill$\square$

For completeness, an explicit finite CW complex shows why this conclusion is too strong for $W_\infty$.  Let $K$ have one vertex, three oriented edges $x,y,z$, and three two-cells with relators
\[
 x^2(xyz)^{-1},\qquad y^3(xyz)^{-1},\qquad z^5(xyz)^{-1}.
\]
Its cellular boundary $C_2(K;\mathbb Z)\to C_1(K;\mathbb Z)$ has, in these bases, the matrix
\[
 M=\begin{pmatrix}1&-1&-1\\-1&2&-1\\-1&-1&4\end{pmatrix},
 \qquad \det M=-1.
\]
It is an isomorphism.  The other reduced cellular groups are zero, so $K$ is acyclic.  Its fundamental group is nontrivial: with permutation multiplication read from right to left, send
\[
 x\longmapsto(1\,2)(3\,4),\quad
 y\longmapsto(1\,3\,5),\quad
 z\longmapsto(1\,2\,5\,3\,4).
\]
These permutations have orders $2,3,5$ respectively and product $xyz=1$, so they satisfy the three relators and define a homomorphism with nontrivial image in $A_5$.  A finite simplicial complex homotopy equivalent to $K$ exists by \citet[Theorem~2C.5]{Hatcher2002}; let $A$ be its poset of nonempty faces.  The nerve of $A$ is its barycentric subdivision.  Thus $BA$ is connected and acyclic but has nontrivial fundamental group.

Apply the construction in the proof to this finite poset.  Then $C_1$, $C_2$, and $C$ are weakly contractible, while $A$ is not.  Hence $C_2\to C$ belongs to $W_\infty$ but $A\to C_1$ does not.  This disproves full W(7 bis) for $W_\infty$.  Mapping this closed-cover square to the trivial square of one-point categories also disproves the implication $b\mathbin{\&}c\Rightarrow a$ in full W(7).  The same example disproves the corresponding converses in Proposition~13 and the corollaries if their extra hypotheses are deleted.  It does not contradict the propositions as actually stated.

The topology behind the example is classical; compare the acyclic presentation complex in \citet[Example~2.38]{Hatcher2002}.  We checked that example, the manuscript's definitions of W(7 bis) and W(7), and the standard nerve and suspension arguments.  We did not find this exact formulation as a test of these manuscript axioms in the references inspected.  The proposition records the deduction needed here, without a claim of priority.  In particular the assertion earlier in the chapter that $W_\infty$ satisfies full W(7) cannot be justified by a Mayer--Vietoris argument alone: local systems on the contractible union restrict to constant systems on $A$, and do not test its nontrivial fundamental group.

There is a precise positive form of the five-lemma intuition.  For a closed cover, every nondegenerate nerve simplex belongs to one of the two pieces, so normalized integral chains form a natural short exact sequence
\[
 0\longrightarrow C_*(NA)
  \longrightarrow C_*(NC_1)\oplus C_*(NC_2)
  \longrightarrow C_*(NC)\longrightarrow0.
 \tag{A7}
\]
The first map has signs $(+,-)$ and the second adds the two inclusions.  Exactness follows separately on each simplex, including degree zero.  Given a map of two covers, if any two of the maps on the three terms of (A7) are quasi-isomorphisms, the third is a quasi-isomorphism by the associated long exact homology sequences.  For the middle term this is equivalent to being a homology isomorphism on each piece.  Thus all three implications of the cube axiom hold for integral homology equivalences of closed covers.  The passage from this homological statement to weak homotopy equivalences is exactly the passage that the counterexample prevents.
\end{dcomment}


\Needspace{13\baselineskip}
\textbf{Corollaire :} \textit{Considérons} deux \textit{diagrammes de type} $\Phi^{\circ}$\textit{, et un homomorphisme} $f = (f_0, f_1, f_2)$ \textit{de l'un dans l'autre (cf. diagramme en traits pleins ci-contre} \texttt{[ci-dessous]}\textit{):}
\[
\begin{tikzcd}[column sep=small, row sep=small, ampersand replacement=\&]
 \& X_1 \arrow[dd, "f_1"'] \arrow[drr, dotted, "i_1"] \& \& \\
X_0 \arrow[ur, "\alpha_1"] \arrow[drr, "\alpha_2"] \arrow[dd, "f_0"'] \& \& \& \cX \arrow[dd, dotted, "f_{\cX}"] \\
 \& X'_1 \arrow[drr, dotted, "i'_1"] \& X_2 \arrow[ur, dotted, "i_2"] \arrow[dd, "f_2"'] \& \\
X'_0 \arrow[ur, "\alpha'_1"] \arrow[drr, "\alpha'_2"] \& \& \& \cX', \\
 \& \& X'_2 \arrow[ur, dotted, "i'_2"] \&
\end{tikzcd}
\]
\textit{d'où un homomorphisme}
\msp{page 127}
\[
f_{\cX} = \int f : \quad \cX = \int
\begin{tikzcd}[column sep=small, row sep=tiny, ampersand replacement=\&, baseline=(\tikzcdmatrixname-2-1.base)]
 \& X_1 \\
X_0 \arrow[ur] \arrow[dr] \& \\
 \& X_2
\end{tikzcd}
\qquad\longrightarrow\qquad
\cX' = \int
\begin{tikzcd}[column sep=small, row sep=tiny, ampersand replacement=\&, baseline=(\tikzcdmatrixname-2-1.base)]
 \& X'_1 \\
X'_0 \arrow[ur] \arrow[dr] \& \\
 \& X'_2,
\end{tikzcd}
\]

\printedpage{86}
% head: A. Grothendieck : Dérivateurs, Ch. XII, Édition des Universités de Montpellier II et Paris VII
\textit{s'insérant dans un diagramme cubique (cf. diagramme précédent complété par des pointillés). Considérons les conditions suivantes.}

\begin{itemize}
\item[\textit{a)}] $f_0\in W$\textit{.}
\item[\textit{b)}] $f_1$ \textit{et} $f_2$ \textit{sont} \texttt{[dans]} $W$\textit{.}
\item[\textit{c)}] $f_{\cX} = \int(f_0, f_1, f_2)\in W$\textit{.}
\end{itemize}

\textit{Supposons l'axiome} W(7) (`lemme des 5', cf. page 76) \textit{satisfait. Alors deux parmi les conditions a), b), c) impliquent la troisième}
\[
\left\{
\begin{aligned}
a \mathbin{\&} b &\quad\Longrightarrow\quad c\\
b \mathbin{\&} c &\quad\Longrightarrow\quad a\\
c \mathbin{\&} a &\quad\Longrightarrow\quad b.
\end{aligned}
\right.
\]

En effet, on a un homomorphisme du diagramme cubique ci-dessus dans le diagramme similaire, où $X_1$, $X_2$, $X'_1$, $X'_2$ sont remplacés respectivement pas les fermés % [sic] pas -> par, misprint
$1\backslash\cX$, $2\backslash\cX$ de $\cX$, et $1\backslash\cX'$, $2\backslash\cX'$ de $\cX'$, les deux carrés supérieur et inférieur étant des inclusions d'ouverts \texttt{[plutôt de fermés]} de $\cX$ resp. de $\cX'$. Ainsi, les conditions a), b), c) sont équivalentes respectivement aux conditions similaires pour le nouveau diagramme cubique
\msp{page 128}
(a) et c) n'ayant d'ailleurs pas changé, et b) donnant
\begin{itemize}
\item[b$'$)] $f'_1 : \cX/1 \longrightarrow \cX'/1$\ \ et\ \ $f'_2 : \cX/2 \longrightarrow \cX'/2$\ \ sont dans $W$.
\end{itemize}
\texttt{[plutôt} $f'_1 : 1\backslash\cX \longrightarrow 1\backslash\cX'$ \texttt{et} $f'_2 : 2\backslash\cX \longrightarrow 2\backslash\cX'$\texttt{]}). Alors la conclusion du corollaire est équivalente à W(7) (formulation duale) dans le cas du nouveau diagramme cubique, q.e.d.

\subsection*{3.2\quad Comparaison de \ $\displaystyle\int
\begin{tikzcd}[column sep=small, row sep=tiny, ampersand replacement=\&, baseline=(\tikzcdmatrixname-2-1.base)]
 \& X_1 \\
X_0 \arrow[ur] \arrow[dr] \& \\
 \& X_2
\end{tikzcd}$ \ avec une somme amalgamée ordinaire.\\
Carrés $W$-homotopiquement cocartésiens.}

La catégorie intégrée $\cX = \displaystyle\int
\begin{tikzcd}[column sep=small, row sep=tiny, ampersand replacement=\&, baseline=(\tikzcdmatrixname-2-1.base)]
 \& X_1 \\
X_0 \arrow[ur] \arrow[dr] \& \\
 \& X_2
\end{tikzcd}$ est solution du problème universel suivant, variante 2-catégorique de celui qui définit la limite inductive des diagrammes. Il s'agit, pour $Y$ objet variable dans Cat, de considérer les systèmes de trois foncteurs $f_0$, $f_1$, $f_2$ et des deux homomorphismes fonctoriels $\lambda_1$, $\lambda_2$ :
\printedpage{87}
% head: A. Grothendieck : Dérivateurs, Ch. XII, Édition des Universités de Montpellier II et Paris VII
\[
\begin{tikzcd}[column sep=large, row sep=large, ampersand replacement=\&]
 \& X_0 \arrow[dl, "\alpha_1"'] \arrow[dr, "\alpha_2"] \arrow[dd, "f_0", ""{name=M, pos=0.62}] \& \\
X_1 \arrow[dr, "f_1"', ""{name=A, pos=0.45}] \& \& X_2 \arrow[dl, "f_2", ""{name=B, pos=0.45}] \\
 \& Y, \&
\arrow[from=A, to=M, bend left, shorten=3pt, "\lambda_1"]
\arrow[from=B, to=M, bend right, shorten=3pt, "\lambda_2"']
\end{tikzcd}
\tag*{$(*)$}
\]
\[
\begin{aligned}
\lambda_1 :\quad & f_1\alpha_1 \quad\longrightarrow\quad f_0\\
\lambda_2 :\quad & f_2\alpha_2 \quad\longrightarrow\quad f_0.
\end{aligned}
\]

\msp{page 129}
Se donner une flèche de Cat
\[
\cX \longrightarrow Y
\]
`revient au même' que de se donner un système $(f_0, f_1, f_2, \lambda_1, \lambda_2)$ comme dessus, et ceci `fonctoriellement en $Y$'. C'est là le cas particulier de la détermination générale des foncteurs de la catégorie fibrée `intégrée' d'un contrafoncteur
\[
I^{\circ} \longrightarrow \Cat, \qquad
\begin{aligned}
i &\quad\longmapsto\quad X_i\\
u &\quad\longmapsto\quad u^*
\end{aligned}
\]
vers une catégorie $Y$, en termes d'un système de foncteurs
\[
X_i \xrightarrow{f_i} Y,
\]
et d'homomorphismes fonctoriels
\[
f_i\circ u^* \xrightarrow{\lambda_u} f_j
\]
pour $u : i \longrightarrow j$, d'où $u^* : X_j \longrightarrow X_i$, avec des conditions de transitivité pour des flèches composées, dont on n'a pas à tenir compte dans le cas actuel, vu qu'il n'y a pas de flèches composées à regarder dans $I = \Phi$ (sauf des composées avec une flèche identique, qu'on peut se dispenser de regarder ici \dots).

\dgrblock{dgrb463}{The universal property keeps its two-cells}
\begin{dcomment}{exposition}{p129:vu qu'il n'y a pas de flèches composées à regarder}{The universal property keeps its two-cells}
The universal property on these pages explains why integration is a useful replacement for an ordinary colimit.  It represents functors equipped with comparison transformations.  Those transformations are part of the input and are allowed to be noninvertible.

We can prove the assertion directly from (A4).  Suppose given $f_i:X_i\to Y$ and $\lambda_k:f_k\alpha_k\Rightarrow f_0$ for $k=1,2$.  Define $F:\mathcal X\to Y$ to be $f_i$ on each fibre.  A mixed morphism from $a\in X_k$ to $x\in X_0$, represented by $h:a\to\alpha_kx$, must be sent to
\[
 f_k(a)\xrightarrow{f_k(h)}f_k(\alpha_kx)
       \xrightarrow{\lambda_{k,x}}f_0(x).
 \tag{A8}
\]
This rule respects precomposition in $X_k$.  It respects postcomposition by $v:x\to x'$ because naturality says
\[
 f_0(v)\lambda_{k,x}=\lambda_{k,x'}f_k(\alpha_kv).
\]
There are no composites of two mixed arrows, so these checks prove functoriality.  Conversely, restrict $F$ to the fibres and apply it to the mixed arrows represented by identities.  This recovers the $f_i$ and the $\lambda_k$, and every mixed arrow factors as an internal arrow followed by one of these identities.  The constructions are inverse.

The assertion also holds at the level of categories of functors.  A natural transformation $F\Rightarrow G$ corresponds to transformations $\theta_i:f_i\Rightarrow g_i$ such that, for $k=1,2$,
\[
 \theta_0\lambda_k=\mu_k(\theta_k\alpha_k),
\]
where $\mu_k$ are the comparison transformations for $G$.  Necessity is naturality on the canonical mixed arrows; sufficiency follows from their factorization just used.  Thus the universal property is a natural isomorphism of categories with the category of these lax cocones.  The word ``lax'' here is fixed by the displayed direction $f_k\alpha_k\to f_0$, so it carries no hidden variance convention.

For a general strict functor $I^{\mathrm{op}}\to\Cat$, the additional condition is explicit.  If $u:i\to j$ and $v:j\to k$, with maps $u^*:X_j\to X_i$ and $v^*:X_k\to X_j$, then
\[
 \lambda_{vu}=\lambda_v\circ(\lambda_u v^*),
 \qquad \lambda_{\mathrm{id}_i}=\mathrm{id}_{f_i}.
\]
For the two-legged shape these conditions have no nonidentity instances.  This is why the construction can be given here without a separate coherence argument.

When the original square is strictly commutative, set $f_0=f_k\alpha_k$ and take both $\lambda_k$ to be identities.  Formula (A8) then defines $\varphi_Q$ on every morphism and shows immediately that $\varphi_Qi_i=f_i$.  Every chosen cartesian transition arrow is sent to an identity.  This supplies the requested construction in footnote~87; the comma-category calculations announced there are carried out in the later lemmas on pp.~152--154.  Definition~3 asks whether this particular comparison belongs to $W$.  A bare commutative square in $\Hot_W$ would not determine the comparison, because it has forgotten the transformations from which it was made.
\end{dcomment}


\msp{page 130}
Si en particulier on a un diagramme \textit{commutatif} dans Cat

% ---- from Der-ChXIIb-012_p046-050.tex
\printedpage{88}
% head: A. Grothendieck : Dérivateurs, Ch. XII, Édition des Universités de Montpellier II et Paris VII
\[
Q \quad = \qquad
\begin{tikzcd}
 & X_0 \arrow[dl, "\alpha_1"'] \arrow[dr, "\alpha_2"] & \\
X_1 \arrow[dr, "f_1"'] & & X_2 \arrow[dl, "f_2"] \\
 & Y, &
\end{tikzcd}
\]
on peut le compléter par
\[
f_0 \overset{\text{déf}}{=} f_1\alpha_1 = f_2\alpha_2
\]
en un diagramme
\[
\begin{tikzcd}[row sep=large, column sep=large]
 & X_0 \arrow[dl, "\alpha_1"'] \arrow[dr, "\alpha_2"] \arrow[dd, "f_0", ""'{name=F0, pos=0.65}] & \\
X_1 \arrow[dr, "f_1"', ""{name=F1, pos=0.4}] & & X_2 \arrow[dl, "f_2", ""'{name=F2, pos=0.4}] \\
 & Y, &
\arrow[from=F1, to=F0, bend left, "\mathrm{id}", shorten <=2pt, shorten >=2pt]
\arrow[from=F2, to=F0, bend right, "\mathrm{id}"', shorten <=2pt, shorten >=2pt]
\end{tikzcd}
\]
ce qui montre qu'on a une donnée du type $(*)$ ci-dessus, définissant donc une flèche
\[
\varphi_Q : \quad \sX = \int
\begin{tikzcd}[row sep=small, column sep=small]
 & X_1 \\
X_0 \arrow[ur] \arrow[dr] & \\
 & X_2
\end{tikzcd}
\qquad \longrightarrow \qquad Y
\]
qui donne lieu à un diagramme cubique

\msp{page 131}
\[
\begin{tikzcd}[row sep=small, column sep=small]
 & X_1 \arrow[dd, "\mathrm{id}"'] \arrow[rrd, "i_1", ""'{name=I1}] & & \\
X_0 \arrow[ru, "\alpha_1"] \arrow[rrr, dotted, "i_0", ""'{name=I0, pos=0.8}] \arrow[rrd, "\alpha_2"] \arrow[dd, "\mathrm{id}"'] & & & \sX \arrow[dd, "\varphi"] \\
 & X_1 \arrow[rrd, "f_1"] & X_2 \arrow[ru, "i_2"', ""{name=I2}] \arrow[dd, "\mathrm{id}"'] & \\
X_0 \arrow[ru, "\alpha_1"] \arrow[rrr, dotted, "f_0"] \arrow[rrd, "\alpha_2"] & & & Y, \\
 & & X_2 \arrow[ru, "f_2"'] &
\arrow[from=I1, to=I0, "\lambda_1"', shorten <=1pt, shorten >=1pt]
\arrow[from=I2, to=I0, bend left, "\lambda_2", shorten <=1pt, shorten >=1pt] % Resolved from PDF leaf 11: a solid curve, unlike the genuinely dotted i_0 line.
\end{tikzcd}
\]
% foot: Version 7 décembre 2002

\printedpage{89}
% head: A. Grothendieck : Dérivateurs, Ch. XII, Édition des Universités de Montpellier II et Paris VII
commutatif sauf le carré supérieur, i.e. on a
\[
\varphi i_1 = f_1, \qquad \varphi i_2 = f_2, \qquad \text{et de plus} \quad \varphi i_0 = f_0,
\]
$\varphi$ transformant de plus les homomorphismes fonctoriels
\[
\lambda_1 : i_1\alpha_1 \longrightarrow i_0, \qquad \lambda_2 : i_2\alpha_2 \longrightarrow i_0
\]
en les homomorphismes fonctoriels identiques pertinents.

\setcounter{footnote}{85}
\textbf{Définition 3.} On dit que le diagramme commutatif
\[
\begin{tikzcd}[row sep=small, column sep=small]
 & X_0 \arrow[dl, "\alpha_1"'] \arrow[dr, "\alpha_2"] & \\
X_1 \arrow[dr, "f_1"] & & X_2 \arrow[dl, "f_2"'] \\
 & Y &
\end{tikzcd}
\tag*{$(*)$}
\]
est \textit{$W$-homotopiquement cocartésien} (\footnote{ou simplement $W$-cocartésien, ou qu'il définit $Y$ comme somme amalgamée $W$-homotopique de etc.}) si l'homomorphisme canonique
\[
\sX = \int
\begin{tikzcd}[row sep=small, column sep=small]
 & X_1 \\
X_0 \arrow[ur] \arrow[dr] & \\
 & X_2
\end{tikzcd}
\qquad \xrightarrow{\ \varphi\ } \qquad Y
\]
qu'il définit est dans $W$. (\footnote{Introduire $\varphi$ dans une proposition en forme, et donner ses propriétés essentielles : flèches strictement cartésiennes sont transformées en flèches \textit{identiques} $+$ lemmes 1 et 2 des p. 152, 154.})

\msp{page 132}
On a alors le :

\textbf{Corollaire 1.} \textit{Supposons} W(7 bis)\textit{. Si le carré} $(*)$ \textit{est $W$-homotopiquement cocartésien, alors}
\[
\left\{
\begin{aligned}
\alpha_1 \in W &\quad\iff\quad f_2 \in W \\
\alpha_2 \in W &\quad\iff\quad f_1 \in W.
\end{aligned}
\right.
\]

Cela résulte aussitôt du diagramme cubique de la page précédente, où toutes les flèches verticales sont dans $W$, et de l'énoncé similaire concernant la face supérieure du cube, donné dans la proposition 13, p. 125.

\Needspace{13\baselineskip}
\textbf{Corollaire 2.} \textit{Supposons} W(7)\textit{. Considérons un diagramme cubique}
% foot: Version 7 décembre 2002

\printedpage{90}
% head: A. Grothendieck : Dérivateurs, Ch. XII, Édition des Universités de Montpellier II et Paris VII
\[
\begin{tikzcd}[row sep=small, column sep=small]
 & X_1 \arrow[dd, "p_1"'] \arrow[rrd, "f_1"] & & \\
X_0 \arrow[ru, "\alpha_1"] \arrow[rrd, "\alpha_2"] \arrow[dd, "p_0"'] & & & Y \arrow[dd, "p"] \\
 & X'_1 \arrow[rrd, "f'_1"] & X_2 \arrow[ru, "f_2"] \arrow[dd, "p_2"'] & \\
X'_0 \arrow[ru, "\alpha'_1"] \arrow[rrd, "\alpha'_2"] & & & Y', \\
 & & X'_2 \arrow[ru, "f'_2"] &
\end{tikzcd}
\]
\textit{où les deux faces inférieure et supérieure du cube sont des carrés $W$-homotopiquement cocartésiens. Considérons les conditions}

\msp{page 133}
\begin{enumerate}[label={}]
\item[\textit{a)}] $p_0 \in W$\textit{.}
\item[\textit{b)}] $p_1$ \textit{et} $p_2 \in W$\textit{.}
\item[\textit{c)}] $p \in W$\textit{.}
\end{enumerate}

\textit{Alors deux de ces conditions impliquent la troisième.}

Résulte aussitôt de la définition 3, et du résultat similaire (cor. page 126) pour le diagramme cubique similaire, avec $Y$, $Y'$ remplacées par $\sX = \int
\begin{tikzcd}[row sep=small, column sep=small]
 & X_1 \\
X_0 \arrow[ur] \arrow[dr] & \\
 & X_2
\end{tikzcd}$
et $\sX' = \int
\begin{tikzcd}[row sep=small, column sep=small]
 & X'_1 \\
X'_0 \arrow[ur] \arrow[dr] & \\
 & X'_2
\end{tikzcd}$\;.

\section*{3.3\quad Intégration propre versus intégration lisse. Énoncé d'équivalence entre les deux points de vue. Commentaires sur le cas des limites $W$-homotopiques quelconques.}

La définition 3, et les corollaires précédents, n'ont d'intérêt que dans la mesure où on dispose de critères maniables, permettant d'affirmer qu'un carré commutatif est bel et bien homotopiquement cocartésien. C'est ce qu'on va essayer de faire tantôt. Il y a aussi une perplexité, concernant le rôle privilégié joué, dans la théorie développée ici, par les immersions fermées, à l'exclusion des immersions ouvertes. Nous avons en somme ramené, du
% foot: Version 7 décembre 2002
\printedpage{91}
% head: A. Grothendieck : Dérivateurs, Ch. XII, Édition des Universités de Montpellier II et Paris VII
\msp{page 134}
point de vue homotopique, les diagrammes quelconques $\begin{tikzcd}[row sep=small, column sep=small]
 & X_1 \\
X_0 \arrow[ur] \arrow[dr] & \\
 & X_2
\end{tikzcd}$ dans Cat, au cas d'un diagramme d'immersions fermées, où la somme amalgamée prend des allures familières. Mais pourquoi pas aussi bien le ramener au cas d'immersions ouvertes? On considérerait alors le diagramme de départ comme définissant une catégorie \textit{cofibrée} sur
\[
\Phi^{\circ} = \Psi = \
\begin{tikzcd}[row sep=small, column sep=small]
 & 1 \\
0 \arrow[ur] \arrow[dr] & \\
 & 2,
\end{tikzcd}
\]
soit $\tX$. On a dans $\Psi$ les deux ouverts
\[
\underbrace{\Psi/1}_{=\{0,1\}} = U_1, \qquad \underbrace{\Psi/2}_{=\{0,2\}} = U_2, \qquad U_1 \cap U_2 = \{0\},
\]
d'où dans $\tX$ deux ouverts images inverses
\[
\tX/1, \qquad \tX/2, \qquad \text{avec } \tX/1 \cap \tX/2 = X_0,
\]
d'où un diagramme commutatif d'inclusions ouvertes
\[
\begin{tikzcd}
 & X_0 \arrow[dl, "\alpha'_1"'] \arrow[dr, "\alpha'_2"] & \\
\tX/1 \arrow[dr, "i'_1"'] & & \tX/2 \arrow[dl, "i'_2"] \\
 & \tX, &
\end{tikzcd}
\tag*{$(*)$}
\]

\dgrblock{dgrb464}{The other orientation of integration}
\begin{dcomment}{exposition}{p133:La définition 3, et les corollaires précédents}{The other orientation of integration}
Grothendieck wants the definition of a homotopy pushout to belong to the span itself, without retaining the auxiliary choice of open or closed inclusions used to construct it. The question is quite concrete. The smooth integral introduced above has transition arrows pointing from the two outer fibres towards the central fibre. The proper integral has transition arrows in the other direction. Neither construction is obtained by merely changing the names of the objects.

Write $A=X_0$, $B=X_1$, $C=X_2$, and let $a:A\to B$ and $b:A\to C$ be the given functors. The proper integral $P$ has the objects of $A$, $B$, and $C$, retained as three disjoint sets. In addition to the arrows within those categories, its mixed hom-sets are
\[
 P(x,y)=B(a(x),y),\qquad
 P(x,z)=C(b(x),z)
 \quad(x\in A,\ y\in B,\ z\in C).
\]
There are no mixed arrows in the opposite directions, or between $B$ and $C$. Composition is composition in $B$ and $C$, with $a$ and $b$ applied when necessary. Thus a functor $P\to Y$ consists of functors $f_A,f_B,f_C$ and transformations
\[
 f_A\longrightarrow f_Ba,\qquad f_A\longrightarrow f_Cb.
\]
This is the reverse of the transformation data described in ``A span is remembered by its transition arrows''. For a strictly commuting cocone these transformations are identities, so both constructions have a canonical comparison with $Y$.

The full subcategory $P_1$ on $A$ and $B$ is a sieve in $P$, as is the corresponding $P_2$ on $A$ and $C$. Their intersection is $A$, and their union is $P$. The inclusion $j:B\to P_1$ has a left adjoint $r$: it is the identity on $B$ and sends $x\in A$ to $a(x)$. The identity
\[
 P_1(x,jy)=B(r(x),y)
\]
proves the adjunction, including its mixed case. Consequently $j$ is a weak equivalence for every fundamental localizer. The same argument applies to $C\to P_2$. This explains the parallel with the closed-cover replacement of the smooth integral: the two replacements have the same three vertex homotopy types, but their transition arrows have opposite orientations.

The equivalences asserted on manuscript page 135 inherit the hypotheses of their earlier counterparts. In particular, the reverse implication from an outer inclusion being in $W$ to the corresponding original leg being in $W$ uses the strong form of W(7 bis). The cube statement uses W(7). These are the additional excision assumptions analysed in ``Why converse excision is an additional axiom''. The comparison between the two integrals, which is the object of Proposition 14, does not require either converse-excision principle.
\end{dcomment}


\msp{page 135}
lequel reçoit le diagramme non commutatif
\[
\begin{tikzcd}[row sep=large, column sep=large]
 & X_0 \arrow[dl, "\alpha_1"'] \arrow[dr, "\alpha_2"] \arrow[dd, dotted, "i_0", ""'{name=J0, pos=0.6}] & \\
X_1 \arrow[dr, "i_1"', ""{name=J1, pos=0.4}] & & X_2 \arrow[dl, "i_2", ""'{name=J2, pos=0.4}] \\
 & \tX, &
\arrow[from=J0, to=J1, dotted, bend right, "\tilde{\lambda}_1"', shorten <=2pt, shorten >=2pt]
\arrow[from=J0, to=J2, dotted, bend left, "\tilde{\lambda}_2", shorten <=2pt, shorten >=2pt]
\end{tikzcd}
\]
% foot: Version 7 décembre 2002
\printedpage{92}
% head: A. Grothendieck : Dérivateurs, Ch. XII, Édition des Universités de Montpellier II et Paris VII
les flèches
\[
X_1 \longrightarrow \tX/1, \qquad X_2 \longrightarrow \tX/2
\]
étant dans $W$ (et même $W$-coasphériques - elles ont même un adjoint à gauche).

Donc on trouve, comme précédemment
\[
\left\{
\begin{aligned}
\alpha_1 \in W &\quad\iff\quad i_2 \in W \\
\alpha_2 \in W &\quad\iff\quad i_1 \in W,
\end{aligned}
\right.
\]
et de même on peut formuler un ‘lemme des cinq’ comme le corollaire page 125 \texttt{[plutôt 126]}, mais avec $\tX$, $\tX'$ au lieu de $\sX$, $\sX'$.

De plus, $\tX$ est la solution d'un problème universel dans Cat, en termes de données
\[
\begin{tikzcd}[row sep=large, column sep=large]
 & X_0 \arrow[dl, "\alpha_1"'] \arrow[dr, "\alpha_2"] \arrow[dd, "f_0", ""'{name=G0, pos=0.6}] & \\
X_1 \arrow[dr, "f_1"', ""{name=G1, pos=0.4}] & & X_2 \arrow[dl, "f_2", ""'{name=G2, pos=0.4}] \\
 & Y, &
\arrow[from=G0, to=G1, bend right, "\tilde{\lambda}_1"', shorten <=2pt, shorten >=2pt]
\arrow[from=G0, to=G2, bend left, "\tilde{\lambda}_2", shorten <=2pt, shorten >=2pt]
\end{tikzcd}
\]

\msp{page 136}
la seule différence avec celui formulé à la page 128 étant la direction opposée des flèches $\lambda_1$, $\lambda_2$. Pour tout carré commutatif
\[
\begin{tikzcd}
 & X_0 \arrow[dl, "\alpha_1"'] \arrow[dr, "\alpha_2"] & \\
X_1 \arrow[dr, "f_1"'] & & X_2 \arrow[dl, "f_2"] \\
 & Y, &
\end{tikzcd}
\tag*{$(*)$}
\]
on trouve encore une flèche canonique
\[
\tilde{\varphi} : \tX \longrightarrow Y.
\]

Il serait tout aussi raisonnable de dire que le carré $(*)$ est $W$-homotopiquement cocartésien, lorsque $\tilde{\varphi}$ est \texttt{[dans]} $W$, que d'adopter la définition 3, page 131 - on trouverait exactement les mêmes corollaires. Les deux définitions seraient elles équivalentes?? % [sic] seraient elles -> seraient-elles, hyphen of the inverted subject missing
% foot: Version 7 décembre 2002

% ---- from Der-ChXIIb-012_p051-055.tex
\printedpage{93}
% head: A. Grothendieck : Dérivateurs, Ch. XII, Édition des Universités de Montpellier II et Paris VII
% foot: Version 7 décembre 2002

Je note que $\tsX$ étant cofibrée sur $\Psi=\Phi^\circ$, $\tsX^\circ$ est fibrée sur $\Psi^\circ=\Phi$. De façon précise, on aura
\[
\tsX^\circ = \int
\begin{tikzcd}[row sep=small, column sep=small]
 & (X_1)^\circ \\
(X_0)^\circ \arrow[ur, "\alpha_1^\circ"] \arrow[dr, "\alpha_2^\circ"'] & \\
 & (X_2)^\circ,
\end{tikzcd}
\]
\msp{page 137}
et il s'agit donc de comparer les deux intégrales
\[
\sX = \int
\begin{tikzcd}[row sep=small, column sep=small]
 & X_1 \\
X_0 \arrow[ur, "\alpha_1"] \arrow[dr, "\alpha_2"'] & \\
 & X_2
\end{tikzcd}
, \qquad
\tsX^\circ = \int
\begin{tikzcd}[row sep=small, column sep=small]
 & (X_1)^\circ \\
(X_0)^\circ \arrow[ur, "\alpha_1^\circ"] \arrow[dr, "\alpha_2^\circ"'] & \\
 & (X_2)^\circ
\end{tikzcd}
,
\]
notamment comme objets de $\HOT_W(\Phi)$. Pour cela, je reprends la construction du foncteur
\[
X \longmapsto \Phi(X) : \Cat \longrightarrow \Cat
\]
de I, p. 72-75, mais je préfère ici noter ce foncteur $\Sigma(X)$ (car $\Phi$ désigne déjà une petite catégorie $\begin{tikzcd}[row sep=tiny, column sep=small]
 & 1 \arrow[dl] \\
0 & \\
 & 2 \arrow[ul]
\end{tikzcd}$). On utilisera les homomorphismes fonctoriels en $X$ :
\[
\begin{tikzcd}[column sep=small]
 & \Sigma(X) \arrow[dl, "p_X"'] \arrow[dr, "q_X"] & \\
X & & X^\circ,
\end{tikzcd}
\]
qui sont $W$-asphériques, donc $\in W$. De plus, $\Sigma(X^\circ)=\Sigma(X)$, $p_{X^\circ}=q_X$, $q_{X^\circ}=p_X$. Ainsi, si
\[
\begin{aligned}
X_* \quad &= \qquad
\begin{tikzcd}[ampersand replacement=\&, row sep=small, column sep=small]
 \& X_1 \\
X_0 \arrow[ur, "\alpha_1"] \arrow[dr, "\alpha_2"'] \& \\
 \& X_2
\end{tikzcd}
\qquad \in \quad \Hom(\Phi^\circ,\Cat)\\
X_*^\circ \quad &= \qquad
\begin{tikzcd}[ampersand replacement=\&, row sep=small, column sep=small]
 \& (X_1)^\circ \\
(X_0)^\circ \arrow[ur, "\alpha_1^\circ"] \arrow[dr, "\alpha_2^\circ"'] \& \\
 \& (X_2)^\circ
\end{tikzcd}
\qquad \in \quad \Hom(\Phi^\circ,\Cat),
\end{aligned}
\]
\msp{page 138}
on trouve
\[
\Sigma(X_*) =
\left(
\begin{tikzcd}[row sep=small, column sep=small]
 & \Sigma(X_1) \\
\Sigma(X_0) \arrow[ur, "\Sigma(\alpha_1)"] \arrow[dr, "\Sigma(\alpha_2)"'] & \\
 & \Sigma(X_2)
\end{tikzcd}
\right)
\]
et les flèches
\printedpage{94}
% head: A. Grothendieck : Dérivateurs, Ch. XII, Édition des Universités de Montpellier II et Paris VII
% foot: Version 7 décembre 2002
\[
\begin{tikzcd}[column sep=small]
 & \Sigma(X_*) \arrow[dl, "p_{X_*}"'] \arrow[dr, "q_{X_*}"] & \\
X_* & & X_*^\circ,
\end{tikzcd}
\]
dans $\underline{\Hom}(\Phi^\circ,\Cat)$. Ces flèches sont dans $W$ terme à terme, donc les homomorphismes dans $\Cat$
\[
\begin{tikzcd}[column sep=small]
 & \int\Sigma(X_*) \arrow[dl, "\int p_{X_*}"'] \arrow[dr, "\int q_{X_*}"] & \\
\sX=\int X_* & & \int X_*^\circ=\tsX^\circ
\end{tikzcd}
\]
sont dans $W$.

\dgrblock{dgrb465}{A common resolution of a category and its opposite}
\begin{dcomment}{exposition}{p137}{A common resolution of a category and its opposite}
Grothendieck's comparison uses a resolution that reverses the category without reversing the direction of the functors between categories. This distinction is responsible for the utility of $\Sigma$. Taking opposites is a covariant operation on $\Cat$: a functor $u:A\to B$ gives $u^\circ:A^\circ\to B^\circ$ in the same direction. Applied objectwise, it therefore preserves the shape of the span.

Here is a concrete convention for $\Sigma(C)$. Its objects are arrows $h:x\to y$ of $C$. An arrow from $h:x\to y$ to $h':x'\to y'$ is a pair
\[
 r:x\to x',\qquad s:y'\to y,
 \qquad h=s h'r.
\]
The projections $p_C(h)=x$ and $q_C(h)=y$ take values in $C$ and $C^\circ$, respectively. This is the opposite of one common convention for the twisted-arrow category. Reversing every displayed arrow identifies $\Sigma(C^\circ)$ with $\Sigma(C)$ and interchanges $p$ and $q$. The equality written in the manuscript may be read as this specified canonical identification.

For completeness, the asphericity of the projections can be seen from their fibres. The projection $p_C$ is a split fibration: over $r:x\to x'$ the inverse-image functor sends $h':x'\to y'$ to $h'r:x\to y'$. Its fibre over $x$ is $(x\backslash C)^\circ$, which has the final object $\mathrm{id}_x$. To test asphericity, pull back along $C/c\to C$. The resulting category $\Sigma(C)/c$ is again split fibred, now over $C/c$, with those same aspherical fibres. Since $C/c$ has a final object, the fibre theorem gives that $\Sigma(C)/c$ is $W$-aspherical. Thus $p_C$ is $W$-aspherical. Applying the identical argument to $C^\circ$ proves the assertion for $q_C$.

The importance of the resolution is its strict functoriality. For a diagram $F:I\to\Cat$, it supplies a diagram of natural transformations
\[
 F\ \xleftarrow{p_F}\ \Sigma F\ \xrightarrow{q_F}\ F^\circ
\]
whose components lie in $W$. Integrating the two transformations again gives weak equivalences. At the level of total categories this compares the smooth integral of $F$ with the smooth integral of $F^\circ$, the latter being canonically the opposite of the proper integral of $F$. Applying the same resolution once more to that total category removes the final opposite.

For a strict cocone $Q$ into $Y$, all these maps commute with the comparison maps into $Y$, $\Sigma(Y)$, or $Y^\circ$. The calculation on manuscript page 140 is therefore a calculation in an arrow category, not merely a list of equivalences between their sources. In each commuting square, both horizontal maps belong to $W$; two-out-of-three gives that one vertical comparison belongs to $W$ precisely when the other does. This proves the equivalence of the four conditions in Proposition 14(b), and the same zigzags prove the stronger assertions about coherent diagrams in parts (c) and (d).

This proof is also useful for its scope. It compares the images after inverting the chosen weak equivalences; it does not assert that $C$ and $C^\circ$ are equivalent categories. The fork $0\to1$, $0\to2$ already separates those statements, since it has an initial object and no final object, while its opposite has a final object and no initial object.
\end{dcomment}


Finalement, l'innocent carré commutatif
\[
Q \quad = \qquad
\begin{tikzcd}[row sep=small, column sep=small]
 & X_0 \arrow[dl] \arrow[dr] & \\
X_1 \arrow[dr] & & X_2 \arrow[dl] \\
 & Y &
\end{tikzcd}
\tag*{$(*)$}
\]
(et le diagramme correspondant
\[
\widehat{Q} \quad = \qquad
\begin{tikzcd}[row sep=small, column sep=small]
 & X_0 \arrow[dl] \arrow[dr] & \\
X_1 \arrow[dr] & & X_2 \arrow[dl] \\
 & \sX=\int X_* \arrow[d, "\varphi_Q"] & \\
 & Y &
\end{tikzcd}
\qquad\qquad )
\tag*{$(**)$}
\]
donne naissance, en le transformant
\msp{page 139}
par $X\longmapsto X^\circ$ et par $X\longmapsto\Sigma(X)$, et en tenant compte des homomorphismes fonctoriels en $X$ $\begin{tikzcd}[row sep=small, column sep=small]
 & \Sigma(X) \arrow[dl] \arrow[dr] & \\
X & & X^\circ
\end{tikzcd}$, à trois diagrammes du type $(**)$, i.e. modelés sur la catégorie opposée de
\[
\Gamma \quad = \qquad
\begin{tikzcd}[row sep=small, column sep=small]
 & 0 & \\
1 \arrow[ur, "\alpha_1"] & & 2 \arrow[ul, "\alpha_2"'] \\
 & 3 \arrow[ul, "i_1"] \arrow[ur, "i_2"'] & \\
 & 4 \arrow[u, "\varphi"'] &
\end{tikzcd}
\qquad
\text{\parbox[c]{0.35\textwidth}{\footnotesize (i.e. $\boldsymbol{\Delta}^1\times\boldsymbol{\Delta}^1$ (\footnotemark[88]), auquelle on a rajouté un objet initial)}} % [sic] auquelle -> à laquelle, grammatical form (catégorie is feminine)
\]
\footnotetext[88]{faux, car le carré n'est pas commutatif}
\printedpage{95}
% head: A. Grothendieck : Dérivateurs, Ch. XII, Édition des Universités de Montpellier II et Paris VII
% foot: Version 7 décembre 2002
(\footnote[89]{On a $\alpha_1i_1\neq\alpha_2i_2$, mais la relation $\alpha_1i_1\varphi=\alpha_2i_2\varphi$.}), savoir
\[
\widehat{Q}, \qquad \widehat{\Sigma(Q)}, \qquad \widehat{Q^\circ},
\]
et deux homomorphismes de diagrammes (\texttt{[dans]} $\underline{\Hom}(\Gamma^\circ,\Cat)$)
\[
\begin{tikzcd}[column sep=small]
 & \widehat{\Sigma(Q)} \arrow[dl] \arrow[dr] & \\
\widehat{Q} & & \widehat{Q^\circ}
\end{tikzcd}
\qquad
\text{(déduit de }
\begin{tikzcd}[column sep=small]
 & \Sigma(Q) \arrow[dl] \arrow[dr] & \\
Q & & Q^\circ
\end{tikzcd}
\text{),}
\tag*{$(***)$}
\]
lesquels sont dans $W$ terme à terme (i.e. sont dans $W_\Gamma\subseteq\Fl(\underline{\Hom}(\Gamma^\circ,\Cat))$). Explicitement, en les sommets $i=0,1,2$, resp. $4$ de $\Gamma$, $(***)$ n'est autre que
\[
D_i =
\left(
\begin{tikzcd}[column sep=small]
 & \Sigma(X_i) \arrow[dl, "p_{X_i}"'] \arrow[dr, "q_{X_i}"] & \\
X_i & & X_i^\circ
\end{tikzcd}
\right)
\qquad \text{resp.} \qquad
\begin{tikzcd}[column sep=small]
 & \Sigma(Y) \arrow[dl, "p_Y"'] \arrow[dr, "q_Y"] & \\
Y & & Y^\circ,
\end{tikzcd}
\]
en le sommet 3 c'est
\[
\begin{tikzcd}[column sep=small]
 & \int\Sigma(X_*) \arrow[dl, "\int p_{X_*}"'] \arrow[dr, "\int q_{X_*}"] & \\
\sX=\int X_* & & \int X_*^\circ=\tsX^\circ,
\end{tikzcd}
\]
obtenu par intégration des flèches
\[
\begin{tikzcd}[column sep=small]
 & \Sigma(X_*) \arrow[dl, "p_{X_*}"'] \arrow[dr, "q_{X_*}"] & \\
X_* & & X_*^\circ.
\end{tikzcd}
\]
\msp{page 140}
Regardant la restriction des deux morphismes de diagrammes $(***)$ aux parties des trois diagrammes de type $\Gamma^\circ$ au dessus de la queue $\{3,4\}$ de $\Gamma$ on trouve le diagramme en traits pleins que voici
\[
\begin{tikzcd}[column sep=3em, row sep=normal]
 & \int\Sigma(X_*) \arrow[dl, "W"'] \arrow[dr, "W"] \arrow[dd, "\varphi_{\Sigma(Q)}"'] & & \Sigma\tsX \arrow[dl, dotted, "W"'] \arrow[dr, dotted, "W"] \arrow[dd, dotted, "\Sigma\tilde{\varphi}_Q"] & \\
\sX \arrow[dd, "\varphi_Q"'] & & \tsX^\circ \arrow[dd, "\varphi_{Q^\circ}=(\tilde{\varphi}_Q)^\circ"'] & & \tsX \arrow[dd, dotted, "\tilde{\varphi}_Q"] \\
 & \Sigma(Y) \arrow[dl, "W"'] \arrow[dr, "W"] & & \Sigma(Y) \arrow[dl, dotted, "W"'] \arrow[dr, dotted, "W"] & \\
Y & & Y^\circ & & Y
\end{tikzcd}
\]
\printedpage{96}
% head: A. Grothendieck : Dérivateurs, Ch. XII, Édition des Universités de Montpellier II et Paris VII
% foot: Version 7 décembre 2002
qu j'ai complété par un deuxième volet en pointillés. Dans ce diagramme de quatre carrés accolés, les côtés horizontaux des carrés représentent des flèches de $\Cat$ qui sont dans $W$. On trouve ainsi : % [sic] qu j'ai -> que j'ai, missing letter
\[
\varphi_Q\in W \iff \varphi_{\Sigma Q}\in W \iff (\varphi_{Q^\circ}=\tilde{\varphi}_Q^\circ)\in W \iff \Sigma\tilde{\varphi}_Q\in W \iff \tilde{\varphi}_Q\in W.
\]
On trouve aussi que les cinq flèches verticales sont isomorphes, en tant que flèches de $\Hot_W$. Pour résumer :

\dgrblock{dgrb466}{The indexing category must retain both paths}
\begin{dcomment}{exposition}{p139}{The indexing category must retain both paths}
Grothendieck keeps the entire comparison diagram because the two routes from $X_0$ to the integral become equal only after applying the map to $Y$. The correction in footnotes 88 and 89 is essential to that purpose. Before reaching $Y$, the routes even land in different fibres of the integral. They cannot be made equal by calling the drawing a commutative square.

Let $J=\Gamma^\circ$. A precise presentation has vertices $0,1,2,3,4$, arrows
\[
 0\xrightarrow{a_1}1\xrightarrow{j_1}3\xrightarrow{v}4,
 \qquad
 0\xrightarrow{a_2}2\xrightarrow{j_2}3,
\]
and the single relation $v j_1a_1=v j_2a_2$. The two arrows $j_1a_1,j_2a_2:0\to3$ remain distinct. Thus $4$ is terminal, but the full subcategory on $0,1,2,3$ is the free diamond, not the ordered square $[1]\times[1]$. This is exactly the relation stated by footnote 89, with the orientation reversed.

To include the dotted central inclusion mentioned in footnote 91, add a generator $r:0\to3$ and the relation $vr=vj_1a_1$. The natural transformations connecting $j_1a_1$ and $j_2a_2$ with $r$ are additional two-dimensional data; they are not relations in this ordinary indexing category. Every resolution used in Proposition 14 also preserves $r$, so adding it causes no difficulty for the objectwise weak-equivalence argument. It simply records more of the construction.
\end{dcomment}


\textbf{Proposition 14.} \textit{Considérons un diagramme (carré commutatif)}
\[
Q \quad = \qquad
\begin{tikzcd}[row sep=small, column sep=small]
 & X_0 \arrow[dl] \arrow[dr] & \\
X_1 \arrow[dr] & & X_2 \arrow[dl] \\
 & Y &
\end{tikzcd}
\]
\textit{dans} $\Cat$\textit{, et le diagramme `opposé' déduit par} $X\longmapsto X^\circ$
\[
Q^\circ \quad = \qquad
\begin{tikzcd}[row sep=small, column sep=small]
 & X_0^\circ \arrow[dl] \arrow[dr] & \\
X_1^\circ \arrow[dr] & & X_2^\circ \arrow[dl] \\
 & Y^\circ, &
\end{tikzcd}
\]
\msp{page 141}
\textit{et les flèches correspondantes}
\[
\left\{
\begin{array}{rlclcl}
\varphi_Q : & \sX_Q & = &
\left(\displaystyle\int
\begin{tikzcd}[ampersand replacement=\&, row sep=small, column sep=small]
 \& X_1 \\
X_0 \arrow[ur] \arrow[dr] \& \\
 \& X_2
\end{tikzcd}
\right)
& \longrightarrow & Y \\[2ex]
\varphi_{Q^\circ} : & \sX_{Q^\circ} & = &
\left(\displaystyle\int
\begin{tikzcd}[ampersand replacement=\&, row sep=small, column sep=small]
 \& X_1^\circ \\
X_0^\circ \arrow[ur] \arrow[dr] \& \\
 \& X_2^\circ
\end{tikzcd}
\right)
& \longrightarrow & Y^\circ.
\end{array}
\right.
\]
\textit{Enfin, considérons aussi les flèches $\tilde{\varphi}_Q$, $\tilde{\varphi}_{Q^\circ}$ canoniques, relative à l'intégrale propre $\int_{\mathrm{pr}}$} % [sic] relative -> relatives, agreement with « les flèches »
\printedpage{97}
% head: A. Grothendieck : Dérivateurs, Ch. XII, Édition des Universités de Montpellier II et Paris VII
% foot: Version 7 décembre 2002
\textit{(au lieu de l'intégrale lisse $\int=\int_{\mathrm{lis}}$)}
\[
\left\{
\begin{array}{rlclcl}
\tilde{\varphi}_Q : & \tsX_Q & = &
\left(\displaystyle\int_{\mathrm{pr}}
\begin{tikzcd}[ampersand replacement=\&, row sep=small, column sep=small]
 \& X_1 \\
X_0 \arrow[ur] \arrow[dr] \& \\
 \& X_2
\end{tikzcd}
\right)
& \longrightarrow & Y \\[2ex]
\tilde{\varphi}_{Q^\circ} : & \tsX_{Q^\circ} & = &
\left(\displaystyle\int_{\mathrm{pr}}
\begin{tikzcd}[ampersand replacement=\&, row sep=small, column sep=small]
 \& X_1^\circ \\
X_0^\circ \arrow[ur] \arrow[dr] \& \\
 \& X_2^\circ
\end{tikzcd}
\right)
& \longrightarrow & Y^\circ.
\end{array}
\right.
\]
\textit{On a alors ce qui suit :}

\begin{enumerate}[label={}, leftmargin=2.5em]\itshape
\item[\circled{\textit{a}}] $\varphi_{Q^\circ}\simeq(\tilde{\varphi}_Q)^\circ$, $\varphi_Q\simeq(\tilde{\varphi}_{Q^\circ})^\circ$ (d'où résulte déjà que $\left\{\begin{array}{lcl}\varphi_Q\in W & \Longleftrightarrow & \tilde{\varphi}_{Q^\circ}\in W\\ \varphi_{Q^\circ}\in W & \Longleftrightarrow & \tilde{\varphi}_Q\in W\end{array}\right.$ (\footnote[90]{\upshape et que les flèches $\varphi_Q$ et $\tilde{\varphi}_{Q^\circ}$ d'une part, $\varphi_{Q^\circ}$ et $\tilde{\varphi}_Q$ d'autre part, sont isomorphes comme flèches dans $\Hot_W$.})).

\item[\circled{\textit{b}}] Les relations
\[
\varphi_Q\in W, \qquad \varphi_{Q^\circ}\in W, \qquad \tilde{\varphi}_Q\in W, \qquad \tilde{\varphi}_{Q^\circ}\in W
\]
sont équivalentes, en d'autres termes : pour que le carré soit $W$-homotopiquement cocartésien (version intégrale lisse), il faut et il suffit que le carré opposé $Q^\circ$ le soit, et ces conditions équivalent aussi aux conditions correspondantes en termes des intégrales propres (au lieu des intégrales lisses).

\item[\circled{\textit{c}}] Les quatre flèches $\varphi_Q$, $\varphi_{Q^\circ}$, $\tilde{\varphi}_Q$, $\tilde{\varphi}_{Q^\circ}$ de $\Cat$ en tant que flèches dans $\Hot_W$, sont isomorphes dans $\underline{\Fl}(\Hot_W)$. Plus précisément, en tant que flèches de $\Cat$,
\msp{page 142}
elles définissent des objets isomorphes de $\HOT_W(\boldsymbol{\Delta}^1)$.

\item[\circled{\textit{d}}] Plus précisément encore, les quatre diagrammes dans $\Cat$ de type $\Gamma^\circ$, où
\[
\Gamma = \;
\begin{tikzcd}[row sep=tiny, column sep=small]
 & 1 \arrow[dl, "\alpha_1"'] & & \\
0 & & 3 \arrow[ul, "i_1"'] \arrow[dl, "i_2"] & 4, \arrow[l, "\varphi"'] \\
 & 2 \arrow[ul, "\alpha_2"] & &
\end{tikzcd}
\]
$\alpha_1i_1\varphi=\alpha_2i_2\varphi$, dans lesquels s'insèrent les quatre flèches qu'on vient d'envisager, savoir
\[
\widehat{Q} = 
\begin{tikzcd}[row sep=tiny, column sep=small]
 & X_1 \arrow[dr, "i_1"] & & \\
X_0 \arrow[ur, "\alpha_1"] \arrow[dr, "\alpha_2"'] \arrow[rr, dotted] & & \sX_Q \arrow[r, "\varphi_Q"] & Y, \\
 & X_2 \arrow[ur, "i_2"'] & &
\end{tikzcd}
\]
\end{enumerate}

% ---- from Der-ChXIIb-012_p056-060.tex
\setcounter{footnote}{90}

\printedpage{98}
% head: A. Grothendieck : Dérivateurs, Ch. XII, Édition des Universités de Montpellier II et Paris VII
% foot: Version 7 décembre 2002
\[
\begin{tikzcd}[row sep=small, column sep=small]
 & X_1^\circ \arrow[dr,"i^{\prime}_{1}"] & & \\
\widehat{Q^\circ} = X_0^\circ \arrow[ur,"\alpha_1^\circ"] \arrow[dr,"\alpha_2^\circ"'] \arrow[rr,dotted] & & \mathcal{X}_{Q^\circ} \arrow[r] & Y^\circ, \\
 & X_2^\circ \arrow[ur,"i^{\prime}_{2}"'] & &
\end{tikzcd}
\]
\[
\begin{tikzcd}[row sep=small, column sep=small]
 & X_1 \arrow[dr] & & \\
\widehat{Q}\tld = X_0 \arrow[ur] \arrow[dr] \arrow[rr,dotted] & & \tilde{\mathcal{X}}_{Q} \arrow[r] & Y, \\
 & X_2 \arrow[ur] & &
\end{tikzcd}
\]
\[
\begin{tikzcd}[row sep=small, column sep=small]
 & X_1^\circ \arrow[dr] & & \\
\widehat{Q^\circ}\tld = X_0^\circ \arrow[ur] \arrow[dr] \arrow[rr,dotted] & & \tilde{\mathcal{X}}_{Q^\circ} \arrow[r] & Y^\circ, \\
 & X_2^\circ \arrow[ur] & &
\end{tikzcd}
\]
(\footnote{\textbf{NB} On devrait encore compléter ces diagrammes (non commutatifs) par les flèches en pointillés \dots}) \textit{sont isomorphes dans} $\Hot_W(\Gamma)$\textit{. De façon plus précise encore, on a les isomorphismes dans} $\uHom(\Gamma^\circ,\Cat)$
\[
(\widehat{Q}\tld)^\circ \simeq \widehat{Q^\circ},
\]
\textit{d'où en l'appliquant à $Q^\circ$ au lieu de $Q$,}
\[
(\widehat{Q^\circ}\tld)^\circ \simeq \widehat{Q},
\]
\textit{ce qui implique (via les isomorphismes canoniques valables dans} $\HOT_W(S)$ \textit{pour tout objet $S$ dans} $\Cat$ \textit{et tout} $X_* \in \mathrm{Ob}\,\uHom(S,\Cat)$\textit{,} $X_* \simeq X_*^\circ$ (\footnote{via le foncteur $\Sigma$ utilisé plus haut \dots})\textit{), qu'on a les isomorphismes canoniques}
\[
\left.\begin{aligned}
\widehat{Q} &\simeq \widehat{Q^\circ}\tld \\
\widehat{Q^\circ} &\simeq \widehat{Q}\tld
\end{aligned}\right\}
\qquad \text{\textit{dans} } \HOT_W(\Gamma),
\]
\textit{mais aussi directement}
\[
\widehat{Q} \simeq \widehat{Q^\circ}, \qquad \widehat{Q}\tld \simeq \widehat{Q^\circ}\tld \qquad \text{\textit{dans} } \HOT_W(\Gamma)
\]
\msp{page 143}\textit{(et il reste à vérifier qu'on a ainsi un système} transitif \textit{d'isomorphismes canoniques dans} $\HOT_W(\Gamma)$\textit{, entre les quatre objets envisagés, ce qui doit être essentiellement sorital \dots).}

\dgrblock{dgrb467}{The four comparisons form a coherent system}
\begin{dcomment}{development}{p143:et il reste à vérifier qu'on a ainsi un système}{The four comparisons form a coherent system}
Grothendieck leaves a small but real task after Proposition 14: the canonical comparisons should compose transitively. Without this check, choosing a route through the four constructions could leave an unspecified automorphism. The following argument supplies the check directly in the localized category of diagrams. We did not find this particular four-comparison verification stated separately in the literature consulted; it is a formal completion of the construction on these pages, with no claim of priority.

\textbf{Coherence of the four integration comparisons.}\ours\
Let $Q$ be a strictly commuting cocone on a small category-valued diagram, and use the same expanded indexing category $J$ for the integral, the cocone, and their structural inclusions. In the localization
\[
 \mathcal L_J=\operatorname{Fun}(J,\Cat)[W_J^{-1}],
\]
the comparisons furnished by $\Sigma$ between the smooth and proper integrals of $Q$ and of $Q^\circ$ form a transitive system of isomorphisms. The assertion includes the arrows from the integral to the cocone vertex and the optional central inclusions.

\emph{Proof.} Write $O$ for objectwise passage to opposite categories. It preserves $W_J$, so it acts on $\mathcal L_J$. The resolution of the preceding block gives a natural isomorphism
\[
 \eta_D:D\xrightarrow{\sim}OD,
 \qquad \eta_D=q_Dp_D^{-1}.
\]
Because $\Sigma(OD)$ is identified with $\Sigma(D)$ with its projections exchanged,
\[
 \eta_{OD}=\eta_D^{-1}.
\]
Naturality holds first for actual maps of diagrams and then for their inverses in the localization; hence it holds for every arrow in $\mathcal L_J$. Applying it to $\eta_D$ gives
\[
 O(\eta_D)\eta_D=\eta_{OD}\eta_D,
 \qquad O(\eta_D)=\eta_D^{-1}.
\]
Thus the opposite operation and its canonical comparison already satisfy the required involution identity.

Now put $A=\widehat Q_{\mathrm{lis}}$ and $B=\widehat{Q^\circ}_{\mathrm{lis}}$. Integrating the objectwise resolution of the original diagram, while using $\Sigma(Y)$ at the cocone vertex, gives
\[
 \alpha_Q:A\xrightarrow{\sim}B.
\]
The same resolution for $Q^\circ$ interchanges its two projections, including at the integrated vertex. Consequently $\alpha_{Q^\circ}=\alpha_Q^{-1}$. By the strict descriptions of the two Grothendieck constructions, the other two objects are canonically
\[
 C=\widehat Q_{\mathrm{pr}}=OB,
 \qquad D=\widehat{Q^\circ}_{\mathrm{pr}}=OA.
\]
Proper integration is the composite of objectwise opposite, smooth integration, and objectwise opposite again. Since $O(\eta_Q)=\eta_Q^{-1}$ also holds in the localized category of input cocones, applying this composite to the resolution shows that its direct proper comparison from $C$ to $D$ is $O(\alpha_Q^{-1})$. Naturality of $\eta$ for $\alpha_Q$ says exactly
\[
 \eta_B\alpha_Q=O(\alpha_Q)\eta_A,
\]
or, equivalently,
\[
 O(\alpha_Q^{-1})\eta_B\alpha_Q=\eta_A.
\]
This is the commutativity of the square formed by the four constructions. Together with the involution identities, it shows that the composite around the square, in either direction, is the identity. Every route between two vertices can be reduced by cancelling inverse adjacent edges and using this square identity. The resulting comparison is independent of the route, and its composites are transitive. All equalities take place in $\mathcal L_J$, so they also hold after restricting to the comparison arrow or to any other subdiagram.\hfill$\square$

The argument is unchanged for an arbitrary index category. For $F:I\to\Cat$, one may define $J$ by freely adjoining a vertex $v$ and arrows $r_i:i\to v$, then a terminal vertex $y$, so that the composites $i\to v\to y$ satisfy the relations of the given strict cocone. The arrows $r_j u$ and $r_i$ are not identified for $u:i\to j$; their images in the integral are linked by transition transformations. With this presentation, the arbitrary-shape claim on manuscript pages 143--144 follows from the same proof. The phrase ``without commutation relations'' describes the first, free adjunction of $v$; only the subsequent vertex $y$ is terminal. Keeping these two operations distinct prevents the indexing error already corrected on page 139.
\end{dcomment}


\subsection*{Remarques.}

\begin{enumerate}
\item[1)] Ce qu'on vient de dire dans le cas des diagrammes dans Cat de type $\Phi^\circ = \Psi$, a l'air de garder un sens pour les diagrammes de type quelconque
\[
X_* : S^\circ \longrightarrow \Cat, \qquad s \longmapsto X_s
\]
\printedpage{99}
% head: A. Grothendieck : Dérivateurs, Ch. XII, Édition des Universités de Montpellier II et Paris VII
% foot: Version 7 décembre 2002
($S$ un objet de Cat). On peut l'intégrer de façon lisse, $\int X_* = \int_{\mathrm{lis}} X_*$, obtenant une catégorie fibrée (scindée) sur $S$, ou de façon propre, $\int_{\mathrm{pr}} X_*$, obtenant une catégorie cofibrée scindée sur $S^\circ$. On peut également considérer $X_*^\circ$, composé de $X_*$ avec l'automorphisme $Y \longmapsto Y^\circ$ dans Cat, et l'intégrer de façon soit propre, soit lisse. Dans tous ces cas, on trouve un diagramme de type $\widehat{S}^\circ$, où $\widehat{S}$ est déduit de $S$ en ajoutant un élément qui s'envoie par une flèche ‘canonique’ (donnée) dans chacun des autres (\textit{sans} relations de commutation!) Un
\msp{page 144}\textit{système inductif} de morphismes dans Cat
\[
f \qquad X_s \xrightarrow{f_s} Y
\]
donne alors un morphisme
\[
\int_{\mathrm{lis}} X_* \xrightarrow{\varphi_Q} Y,
\]
où $Q$ désigne le diagramme, de type $\widehat{S}^\circ$, formé des $X_s$, de $Y$, des $u^* : X_t \longrightarrow X_s$ (pour $u : s \longrightarrow t$ dans $S$) et des $f_s$. De même on aura
\[
\int_{\mathrm{prop}} X_* \xrightarrow{\tilde{\varphi}_Q} Y,
\]
et si
\[
\Gamma = (\widehat{S} + \text{objet initial}),
\]
on trouve des diagrammes $\widehat{Q}$, $\widehat{Q}\tld$ de type $\Gamma^\circ$ dans Cat. Sauf erreur, les relations entre $\widehat{Q}$, $\widehat{Q}\tld$, $\widehat{Q^\circ}$, $\widehat{Q^\circ}\tld$ sont les mêmes que celles explicitées dans la proposition 14 ci-dessus. Je ne vais pas m'arrêter là-dessus avec le soin requis, sauf en cas de besoin. On dira qu'un diagramme $Q$ définit $Y$ comme $\varinjlim\limits_{I}$ $W$-homotopique \texttt{[où} $I = S^\circ$\texttt{]} si $\varphi_Q \in W$. Il revient au même de dire que $\varphi_{Q^\circ} \in W$, ou que $\tilde{\varphi}_Q \in W$ ou que $\tilde{\varphi}_{Q^\circ} \in W$ \dots
\msp{page 145}

\dgrblock{dgrb468}{The comparison map in later formalisms}
\begin{dcomment}{later}{p144}{The comparison map in later formalisms}
The definition by a specified map from an integral has a direct later interpretation. For ordinary homotopy types, put
\[
 \mathbb D(J)=\operatorname{Fun}(J,\Cat)[W_{\infty,J}^{-1}]
             =\HOT_{W_\infty}(J^{\mathrm{op}}).
\]
The opposite in the last expression corrects for the manuscript's presheaf convention. A square is an object of $\mathbb D(\square)$, where $\square=[1]\times[1]$. Let $i:\ulcorner\hookrightarrow\square$ omit the terminal vertex. In the derivator formulation a square $Q$ is cocartesian when the counit
\[
                   i_!i^*Q\longrightarrow Q
\]
is an isomorphism \citep[Definition~3.9]{Groth2013}. At the three vertices of $\ulcorner$ its components are already isomorphisms. At the last vertex, the component is the homotopy-colimit comparison represented here by $\varphi_Q$. This identifies Definition~3 with the later coherent-square definition. It also explains why a square merely drawn in $\mathbb D(e)$ is insufficient: it need not specify the coherent object to which $i_!$ applies.

The composition and cancellation theorem for adjacent cocartesian squares is \citet[Proposition~3.13]{Groth2013}. Its hypotheses concern squares inside one coherent rectangle. They do not assert that an arbitrary map between homotopy pushouts reflects an equivalence on the intersection from equivalences on the pieces and the union. The latter is the stronger cube assertion tested by the acyclic example above.

There is a second distinction when the fibres are retained as categories. A Cartesian fibration classified by a diagram of $\infty$-categories models its colimit only after its Cartesian transition edges have been inverted; see \citet[Corollary~3.3.4.3 and Warning~3.3.4.4]{Lurie2009HTT}. Passing further to the groupoid completion inverts every edge. The corresponding space-valued statement is Proposition~3.3.4.5 of the same reference. This is compatible with the classifying-space interpretation of the manuscript's integral, while preserving the stronger categorical information before localization.

The constant terminal-category diagram over $[1]$ makes the distinction concrete. Its ordinary total category is $[1]$. The colimit of the diagram in categories is $e$, and its colimit in $\mathrm{Cat}_\infty$ is also $e$. The unique nonidentity arrow of $[1]$ is a transition arrow; inverting it gives a category equivalent to $e$. Before that inversion, $[1]$ and $e$ are not equivalent categories: the two objects of $[1]$ are not isomorphic. Their classifying spaces are nevertheless both contractible. Thus the three assertions---a total category, a categorical colimit, and a colimit of homotopy types---agree only after the indicated operations. The present text is computing the last of them and retaining a strict representative for its map to the proposed target.
\end{dcomment}


\item[2)] Le cas ‘universel’, pour
\begin{tikzcd}[baseline=(\tikzcdmatrixname-2-1.base), row sep=small, column sep=small]
 & X_1 \\
X_0 \arrow[ur] \arrow[dr] & \\
 & X_2
\end{tikzcd}\par\smallskip
% Reader layout: separate the two inline diagrams without changing the wording.
donné, d'un carré commutatif $Q = {}$%
\begin{tikzcd}[baseline=(\tikzcdmatrixname-2-1.base), row sep=tiny, column sep=small]
 & X_1 \arrow[dr] & \\
X_0 \arrow[ur] \arrow[dr] & & Y \\
 & X_2 \arrow[ur] &
\end{tikzcd}
dans Cat, donnant naissance à
\[
\begin{tikzcd}[row sep=small, column sep=small]
 & & & X_1 & & & \\
\varphi_Q : &[1em] \left(\vrule height 6ex depth 6ex width 0pt\right. &[-1.5em] \int X_0 \arrow[ur] \arrow[dr] & & \left.\vrule height 6ex depth 6ex width 0pt\right) &[2em] \longrightarrow &[2em] Y, \\
 & & & X_2 & & &
\end{tikzcd}
\]
est celui de $Y = X_1 \amalg_{X_0} X_2$, donnant naissace à % [sic] naissace -> naissance, letter n omitted
\printedpage{100}
% head: A. Grothendieck : Dérivateurs, Ch. XII, Édition des Universités de Montpellier II et Paris VII
% foot: Version 7 décembre 2002
\[
\begin{tikzcd}[row sep=small, column sep=small]
 & X_1 \arrow[dr] & \\
Q_0 = Q(X_*) = X_0 \arrow[ur] \arrow[dr] & & X_1 \amalg_{X_0} X_2. \\
 & X_2 \arrow[ur] &
\end{tikzcd}
\]

Le cas général s'y ramène, car on aura un diagramme commutatif
\[
\begin{tikzcd}[row sep=small, column sep=small]
 & & X_1 & & & \\
\left(\vrule height 6ex depth 6ex width 0pt\right. &[-1.5em] \int X_0 \arrow[ur] \arrow[dr] & & \left.\vrule height 6ex depth 6ex width 0pt\right) \arrow[rr,"{\varphi_{Q_0}=\varphi_{Q(X_*)}}"] \arrow[ddrr,"\varphi_Q"] &[4em] \hphantom{MMMM} &[4em] X_1 \amalg_{X_0} X_2 \arrow[dd,"\alpha_Q"] \\
 & & X_2 & & & \\
 & & & & & Y.
\end{tikzcd}
\]
\end{enumerate}

Mais il doit y avoir des cas intéressants où $\varphi_{Q_0}$ n'est \textit{pas} dans $W$, i.e. le carré Cat-cocartésien
\begin{tikzcd}[baseline=(\tikzcdmatrixname-2-1.base), row sep=tiny, column sep=small]
 & X_1 \arrow[dr] & \\
X_0 \arrow[ur] \arrow[dr] & & X_1 \amalg_{X_0} X_2 \\
 & X_2 \arrow[ur] &
\end{tikzcd}
n'est pas $W$-homotopiquement cocartésien, et où pourtant on peut trouver un carré commutatif $Q$
\msp{page 146}complétant $X_*$, de telle façon que ce carré soit $W$-homotopiquement commutatif \texttt{[plutôt cocartésien]}, i.e. que $\varphi_Q = \alpha_Q \varphi_{Q_0}$ soit dans $W$. C'est d'autant plus plausible que dans Cat, les sommes amalgamées ont tendance à ne pas avoir des bonnes propriétés homotopiques, même si l'une des flèches $\alpha_1$, $\alpha_2$ est \texttt{[un]} monomorphisme (voire une immersion ouverte), contrairement à ce qui se passe dans $\boldsymbol{\Delta}^\wedge$, et plus généralement dans les catégories de la forme $Z^\wedge$ ($Z$ dans Cat).

\dgrblock{dgrb469}{A strict cocone can repair the wrong pushout}
\begin{dcomment}{development}{p146:de telle façon que ce carré soit $W$-homotopiquement}{A strict cocone can repair the wrong pushout}
Grothendieck asks for more than a strict pushout that fails to be a homotopy pushout. He asks whether its failure can disappear after mapping that pushout to a different target, while retaining the original span. The following finite example answers his question affirmatively for $W_\infty$, even when one leg is an open immersion. The sources checked on preservation of pushouts, including the Dwyer-map criterion and its refinement in \citet[Theorems~1.5--1.6, arXiv v3]{HackneyOzornovaRiehlRovelli2024}, do not supply this particular construction. We give a complete argument rather than a priority claim.

\textbf{A finite repair example.}\ours\
There exist finite posets $E\subset B$, with $E$ a sieve in $B$, such that $BE$ and $BB$ are contractible, while the categorical pushout $P=B\amalg_E e$ has classifying space $S^3$. Consequently the canonical comparison
\[
 c:\int_{\mathrm{lis}}(B\leftarrow E\to e)\longrightarrow P
\]
is not in $W_\infty$, but its composite with $P\to e$ is in $W_\infty$. Here $BD=|ND|$ denotes the classifying space of a category $D$.

\emph{Proof.} We first construct a map from a contractible face poset to a spherical one whose comma categories are connected. Let $K$ be the finite simplicial complex with six vertices $a,b,c,A,B,C$ and maximal faces
\[
 abc,\quad abB,\quad aBA,\quad bcC,\quad bCB,\quad caA,\quad cAC.
\]
It is the disk obtained by placing the triangle $abc$ inside the triangle $ABC$ and triangulating the surrounding annulus with the remaining six triangles. In particular it is contractible. This description can also be checked without a drawing: successive free-face collapses remove
\[
\begin{gathered}
(AB,aAB),\ (AC,cAC),\ (aA,acA),\ (A,cA),\ (BC,bBC),\\
(aB,abB),\ (B,bB),\ (bC,bcC),\ (C,cC),\\
(ab,abc),\ (a,ac),\ (b,bc),
\end{gathered}
\]
leaving the vertex $c$. In each pair the first simplex is a free codimension-one face of the second at the moment it is removed.

Let $L$ be the boundary of the tetrahedron on vertices $a,b,c,n$. The vertex map fixing $a,b,c$ and sending $A,B,C$ to $n$ is simplicial from $K$ to $L$: the seven triangles have images $abc$, $abn$, $an$, $bcn$, $bn$, $can$, $cn$. Let $E$ and $S$ be the posets of nonempty faces of $K$ and $L$, and let $f:E\to S$ send a face to its image. Thus $BE$ is a subdivided disk and $BS$ a subdivided two-sphere.

For every face $s$ of $L$, the comma category $E/s$ is the face poset of the subcomplex consisting of faces whose image is contained in $s$. If $n\notin s$, this subcomplex is the simplex on $s$, so it is connected and nonempty. If $n\in s$, it contains the entire outer boundary $AB\cup BC\cup CA$. Every inner vertex allowed by $s$ attaches to that boundary by a radial edge: $a$ by $aA$, $b$ by $bB$, and $c$ by $cC$. Every remaining simplex contains vertices of this same connected subcomplex. It follows that
\[
 \pi_0(E/s)=\{*\}\quad\text{for every }s\in S.
\]
Only connectedness is claimed: for $s=\{n\}$ the comma category is the face poset of the outer boundary circle.

Adjoin an initial object $o$ to $S$ and call the resulting poset $T$. Define a poset $B$ on the disjoint union of $E$ and $T$. Keep their internal orders and add
\[
 e\leq s\quad\Longleftrightarrow\quad f(e)\leq s
 \qquad(e\in E,\ s\in S).
\]
There are no relations from $T$ to $E$, and none from $E$ to $o$. Transitivity follows from monotonicity of $f$. The subposet $E$ is a sieve. The functor $r:B\to T$ which is the identity on $T$ and sends $e$ to $f(e)$ is left adjoint to the inclusion $j:T\to B$, because
\[
 B(b,jt)=T(r(b),t)
\]
for every $b\in B$, $t\in T$, including the empty hom-sets when $t=o$ and $b\in E$. Hence $Bj$ is a homotopy equivalence. The initial object $o$ makes $BT$, and therefore $BB$, contractible.

It remains to compute the categorical pushout, since a topological quotient would give the wrong answer. Collapsing the sieve $E$ to one object $*$ introduces no arrows from $T$ to $*$ and does not change arrows within $T$. Any arrow from $*$ to $s\in S$ is represented by a relation $e\leq s$ in $B$. The relations after collapsing $E$ identify two representatives exactly along zigzags of comparable elements in $E/s$. Thus
\[
 \operatorname{Hom}_{P}(*,s)=\pi_0(E/s)=\{*\},
 \qquad \operatorname{Hom}_{P}(*,o)=\emptyset.
\]
For clarity, this description also verifies the universal property: a functor out of $B$ that is constant on $E$ assigns the same arrow to any two representatives in one such connected component, and therefore factors uniquely through the displayed category. Conversely every functor from that category supplies exactly such compatible data on $B$ and $e$. Composition is induced by the maps $E/s\to E/s'$ for $s\leq s'$.

Consequently $P$ is the poset with two incomparable minimal objects $o,*$, both lying below every object of $S$. Its order complex is exactly the simplicial join
\[
 NP=N\{o,*\}*NS,
 \qquad BP\cong S^0*|\operatorname{sd}L|\cong S^3.
\]
The three categories in the original span $B\leftarrow E\to e$ are contractible. Their smooth integral is a split fibration over the contractible indexing fork, with contractible fibres, and hence is contractible. Its canonical map $c$ cannot be a weak equivalence to $P$, whose third homology is $\mathbb Z$. On the other hand the unique functor from the integral to $e$ is a weak equivalence, and it is precisely the composite $P\to e$ after $c$. This proves the assertion.\hfill$\square$

The construction has $25$ objects in $E$, $14$ in $S$, $40$ in $B$, and $16$ in $P$. The accompanying finite computation checks the face lists, the elementary collapses, all fourteen connected comma categories, and the homology of $NP$; the join calculation above identifies its entire homotopy type. The mechanism is visible in the proof: the ordinary categorical pushout keeps only $\pi_0(E/s)$ in each mixed hom-set. It replaces a family containing a circular comma category by a family of singletons, and the resulting total category acquires the additional sphere. Sending that sphere to a point supplies the strict cocone envisaged on manuscript pages 145--146.
\end{dcomment}


\subsection*{3.4\quad Critères pour les carrés homotopiquement cocartésiens.}

Les critères que j'ai dégagé sont tous contenus dans l'énoncé suivant :

\textbf{Théorème 4} (\footnote{\textbf{NB} Dans ce théorème, et les corollaires 1, 2, 3, on n'a besoin que de la \textit{forme faible} de W(7 bis), savoir $i_0 \in W \Longrightarrow i \in W$ (cf. p. 79) (Loc(6$'$a) avec les notations de XVI).},\footnote{\textbf{NB} Au total, le théorème (partie directe) est valable sans hypothèse W(7 bis) ou W(7 bis faible) sur $W$, et il en est de même de ses corollaires 1 à 3 (à l'exclusion des corollaires 4, 5, qui eux utilisent W(7 bis) dans toute sa force).}). \textit{Considérons un diagramme carré commutatif dans} Cat
\[
Q \quad = \quad
\begin{tikzcd}
X_0 \arrow[r,"\alpha_1"] \arrow[d,"\alpha_2"'] \arrow[dr,dotted,"f_0"] & X_1 \arrow[d,"f_1"] \\
X_2 \arrow[r,"f_2"'] & Y,
\end{tikzcd}
\tag*{$(*)$}
\]
\printedpage{101}
% head: A. Grothendieck : Dérivateurs, Ch. XII, Édition des Universités de Montpellier II et Paris VII
% foot: Version 7 décembre 2002
\textit{et considérons le morphisme canonique}
\[
\begin{tikzcd}[row sep=small, column sep=small]
 & & X_1 & & & \\
\varphi_Q : \quad \mathcal{X} = \Bigl(\int &[-1.2em] X_0 \arrow[ur,"\alpha_1"] \arrow[dr,"\alpha_2"'] & & \Bigr) &[2em] \longrightarrow &[2em] Y, \\
 & & X_2 & & &
\end{tikzcd}
\]
\msp{page 147}\textit{l'intégrale étant l'intégrale} lisse. \textit{(Il y a un énoncé dual concernant l'intégrale propre de}\newline
\begin{tikzcd}[baseline=(\tikzcdmatrixname-2-1.base), row sep=small, column sep=small]
 & X_1 \\
X_0 \arrow[ur] \arrow[dr] & \\
 & X_2
\end{tikzcd}
\textit{.)}

\begin{enumerate}
\item[\textit{a)}] \textit{Pour que $\varphi_Q$ soit $W$-lisse, il faut et il suffit que $f_1$, $f_2$ et $f_0 = f_1\alpha_1 = f_2\alpha_2$ le soient.} (Moyennant W(7 bis) (\footnote{Pour la partie directe ($f_0$, $f_1$, $f_2$ lisses $\Longrightarrow \varphi_Q$ lisse), on n'a besoin que de W(7 bis) (\footnotemark), forme faible. Ce n'est que pour la réciproque (qui ne m'a encore jamais servie) qu'on a besoin de W(7 bis), dans toute sa force.}\footnotetext{On n'a pas même besoin de W(7 bis) faible, cf. annotation page 156. On utilise une hypothèse restrictive sur $W$ dans la réciproque seulement ($\varphi_Q$ lisse $\Longrightarrow f_0$, $f_1$, $f_2$ lisses).})).

\item[\textit{b)}] \textit{La fibre de $\varphi_Q$ en un élément $y$ de $Y$ est isomorphe à}
\[
\begin{tikzcd}[row sep=small, column sep=small]
 & & X_{1\,y} & \\
\int &[-1.2em] X_{0\,y} \arrow[ur] \arrow[dr] & &[2em] (\simeq \mathcal{X}_y). \\
 & & X_{2\,y} &
\end{tikzcd}
\tag*{$(**)$}
\]
\end{enumerate}

\textbf{Corollaire 1.} (\footnote{Valable sans hypothèse sur $W$ (autre que Loc(1,2,3 bis)).}) \textit{Supposons $f_1$, $f_2$ et $f_1\alpha_1 = f_2\alpha_2$ $W$-lisses, et les catégories intégrées $(**)$ $W$-asphériques pour tout $y \in \mathrm{Ob}\,Y$. (NB Il s'agit encore de l'intégrale lisse, mais peu importe au fond, car on a vu (cf. partie supérieure du diagramme p. 140, qui implique que $\mathcal{X}$} \texttt{[est]} \textit{$W$-asphérique si et seulement si $\tilde{X}$ l'est) que l'une est $W$-asphérique si et seulement si l'autre l'est.) Alors $\varphi_Q \in W$, i.e. le carré $(*)$ est $W$-homotopiquement cartésien} \texttt{[plutôt cocartésien]}\textit{. Par dualité, même conclusion si on suppose $f_1$, $f_2$, $f_1\alpha_1 = f_2\alpha_2$} propres \textit{au lieu de lisses. (Car $Q$ $W$-homotopiquement trivial} \texttt{[plutôt cocartésien]}\textit{, si et seulement si $Q^\circ$ l'est, par} \texttt{[la]} \textit{prop. 14.)}

\msp{page 148}En effet, un morphisme dans Cat qui est $W$-lisse et à fibres $W$-asphériques est dans $W$ (et c'est même une $W$-équivalence universelle \dots).

\dgrblock{dgrb470}{Smoothness and the two different uses of a fibre}
\begin{dcomment}{exposition}{p147}{Smoothness and the two different uses of a fibre}
Grothendieck now turns the comparison into a criterion that can be checked on a given square. There are two separate operations in Theorem 4. The fibre formula is an identity of categories: imposing that the comparison map has the value $y$ imposes the same condition in each of its three pieces. Smoothness is a local lifting condition, tested on slices at objects of the integral. The latter condition is what permits a strict fibre to determine the homotopy-theoretic conclusion.

It is useful to read the annotations as a refinement of the hypotheses, rather than as several competing versions of the theorem. If $f_0,f_1,f_2$ are $W$-smooth, then $\varphi_Q$ is $W$-smooth for every fundamental localizer. The reverse implication for $f_0$ is asserted under the strong form of W(7 bis). No extra excision axiom is required for the fibre identity
\[
 (\mathcal X_Q)_y\cong
 \int_{\mathrm{lis}}(X_{1y}\leftarrow X_{0y}\to X_{2y}).
\]
The calculations on manuscript pages 152--157 establish these assertions; their local content is explained below at the two lemmas. In particular, the annotations do not justify removing the strong hypothesis from the converse.

The last step of Corollary 1 deserves to be made explicit. If $u:E\to Y$ is $W$-smooth, the fibre inclusion $E_y\to y\backslash E$ is $W$-aspherical. If every $E_y$ is $W$-aspherical, two-out-of-three gives that every $y\backslash E$ is $W$-aspherical. The dual of the fundamental Theorem A then yields $u\in W$. This proof is stable under arbitrary base change: smoothness is stable under base change and the new strict fibres are the old ones. Hence $u$ is in fact a universal $W$-equivalence. The same argument with properness uses $E/y$ instead. These are the fibre criteria formalized in \citet[Theorem~1.2]{Maltsiniotis2005Asphericite}.

Consequently the forward criterion in Corollary 1 has precisely the strength needed later: smoothness of all three structure maps, together with asphericity of the integrated fibres, proves the square homotopy cocartesian, and even proves the comparison a universal weak equivalence. The qualifications in the first two footnotes to Theorem 4 correctly anticipate this result.
\end{dcomment}


\textbf{Corollaire 2}. \textit{On suppose que $W$ satisfait} W(7 bis) (axiome du carré cocartésien) (\footnote{inutile.}).
\printedpage{102}
% head: A. Grothendieck : Dérivateurs, Ch. XII, Édition des Universités de Montpellier II et Paris VII
% foot: Version 7 décembre 2002
\textit{Supposons le carré $Q$ cocartésien dans} Cat (\footnote{NB On n'utilise pas vraiment que $Q$ soit cocartésien, plutôt que $Q$ est \textit{cartésien,} et $f_2$ (donc aussi $\alpha_1$) une immersion ouverte (resp. fermée), et que pour $y \in Y\backslash X_2$, $X_{1\,y}$ soit $W$-asphérique. (Dans le cas où $Q$ est cocartésien, $X_{1\,y}$ est $\simeq e$.)})\textit{, et que $\alpha_1$ soit une immersion ouverte (resp. fermée) et $f_1$ lisse (resp. propre). Alors $Q$ est $W$-homotopiquement cocartésien.} (\footnote{On a besoin ici de W(7 bis), forme directe, i.e. de Loc(6$'$a) (notation de XVI).})

Il suffit d'établir le résultat non respé, l'autre en résultera par dualité (puisque $Q$ est $W$-cocartésien si et seulement si $Q^\circ$ l'est). Par hypothèse $f_1$ est $W$-lisse, d'autre part $\alpha_1$ et $f_2$ sont des immersions ouvertes, donc $W$-lisses, donc $f_1\alpha_1$ est $W$-lisse aussi. (D'ailleurs $\alpha_2$ aussi, car le carré est aussi cartésien, donc $\alpha_2$ déduit de $f_1$ (qui est $W$-lisse par hypothèse) par changement de base $X_2 \longrightarrow Y$.) Donc $\varphi_Q$ est \texttt{[}$W$\texttt{-]}lisse, et il suffit de voir que ses fibres sont $W$-asphériques. Le calcul $(**)$ donne, pour $y$ dans $Y\backslash X_2$,
\[
\begin{tikzcd}[row sep=small, column sep=small]
 & & e & \\
\mathcal{X}_y \simeq \quad \int &[-1.2em] \emptyset \arrow[ur] \arrow[dr] & &[1em] \simeq e, \qquad \text{O.K.} \\
 & & \emptyset &
\end{tikzcd}
\]
Et pour $y$ dans $X_2$
\[
\begin{tikzcd}[row sep=small, column sep=small]
 & & X_{1\,y} & \\
\mathcal{X}_y \simeq \int &[-1.2em] X_{0\,y} \arrow[ur,"\sim"] \arrow[dr] & &[2em] . \\
 & & e &
\end{tikzcd}
\]
\msp{page 149}Comme $X_{0\,y} \longrightarrow X_{1\,y}$ est un isomorphisme, donc est dans $W$, il s'ensuit que $e \longrightarrow \mathcal{X}_y$ l'est aussi (prop. 13, p. 125), moyennant W(7 bis), donc $\mathcal{X}_y$ est $W$-asphérique, O.K. (\footnote{On n'a besoin que de la forme directe de W(7 bis), i.e. (notation de XVI) de Loc(6$'$a).}).

\textbf{Corollaire 3.} (Moyennant W(7 bis) (\footnote{inutile.}).) \textit{Soit $X$ dans} Cat\textit{, $U_1$, $U_2$ deux ouverts de $X$ tels que $X = U_1 \cup U_2$, $U_0 = U_1 \cap U_2$ l'ouvert intersection. Alors le carré d'immersions ouvertes}
\[
\begin{tikzcd}[row sep=large, column sep=large]
U_0 \arrow[r,hook] \arrow[d,hook] & U_1 \arrow[d,hook] \\
U_2 \arrow[r,hook] & X
\end{tikzcd}
\]
\textit{est $W$-cocartésien. (Énoncé dual pour deux fermés $T_1$, $T_2$ tels que $T_1 \cup T_2 = X$, avec $T_0 = T_1 \cap T_2$.)}

On applique le corollaire 2. On sait que le carré est cocartésien dans Cat (\footnote{On n'a pas même besoin de le savoir, cf. annotation marginale du corollaire 2.}), de plus $U_1 \lhook\joinrel\longrightarrow X$ est bien $W$-lisse, étant une immersion ouverte.

% ---- from Der-ChXIIb-012_p061-065.tex
\printedpage{103}
% head: A. Grothendieck : Dérivateurs, Ch. XII, Édition des Universités de Montpellier II et Paris VII

\noindent\textbf{Corollaire 4.} (Moyennant W(7 bis) (\footnote[104]{Ici on a besoin de W(7 bis) forme forte, i.e. $i_0\in W \iff i\in W$ (cf. annotation th. 4).}).) \textit{Soit $U \hto[i] X$ une immersion ouverte telle que $T=X\backslash U$ soit isomorphe à la catégorie ponctuelle $e$. Pour que $i\in W$, il faut et il suffit que $i$ soit $W$-asphérique, i.e. que $U/t$ (où $t$ est l'unique objet de $T$) soit $W$-asphérique.}

Il suffit bien sûr de prouver que
\[
i\in W \Longrightarrow U/t \ W\text{-asphérique.}
\]

\msp{page 150}
Considérons le carré cartésien
\[
\begin{tikzcd}[row sep=large,column sep=large]
U/t \arrow[r, hook, "i/t"] \arrow[d, "p"'] & X/t \arrow[d, "q"] \\
U \arrow[r, hook, "i"] & X,
\end{tikzcd}
\]
nous allons prouver qu'il est $W$-cocartésien. Moyennant W(7 bis), il s'ensuivra que $i\in W$ implique $i/t\in W$, et comme $X/t$ est $W$-asphérique (ayant un objet final), cela signifie que $U/t$ est $W$-asphérique, ce qu'on veut prouver.

\dgrblock{dgrb471}{Why the cover criterion needs no converse excision}
\begin{dcomment}{examples}{p149}{Why the cover criterion needs no converse excision}
Grothendieck's cover criterion is meant to provide a large supply of actual homotopy pushouts. The annotations correctly remove W(7 bis) from Corollaries 2 and 3, although the displayed proof still invokes it. The missing simplification is elementary: one of the maps in the only nontrivial fibre is an isomorphism, and that special case admits a direct contraction.

For a category $A$, consider
\[
 E=\int_{\mathrm{lis}}(A\xleftarrow{\mathrm{id}}A\longrightarrow e).
\]
Keep the central copy $A_0$ and the outer point $*$, and call their full subcategory $K$. There is a unique arrow $*\to a$ for every $a\in A_0$, so $*$ is initial in $K$. Define $R:E\to K$ to be the identity on $K$ and to send the outer copy $a\in A_1$ to its central copy in $A_0$. On mixed arrows it uses the corresponding arrow of $A$. The canonical transition arrows $a_1\to a_0$ give a natural transformation $\mathrm{id}_E\to jR$, where $j:K\to E$ is the inclusion, and $Rj=\mathrm{id}_K$. Thus $j$ and $R$ are homotopy inverse after taking nerves. Since $K$ has an initial object, $E$ is $W$-aspherical for every fundamental localizer. This proof also covers $A=\emptyset$.

In Corollary 2, a point outside $X_2$ has integrated fibre $e$. At a point of $X_2$, cartesianness identifies $X_{0y}\to X_{1y}$ with an isomorphism, and the fibre is of the form just contracted. The three maps to $Y$ are smooth: $f_1$ by hypothesis, $f_2$ because it is a sieve inclusion, and $f_0=f_1\alpha_1$ by composition. The forward part of Theorem 4 therefore proves the conclusion without W(7 bis). This also establishes the stronger version recorded in footnote 99: strict cocartesianness may be replaced by cartesianness, the sieve condition on $f_2$, smoothness of $f_1$, and $W$-asphericity of $X_{1y}$ outside $X_2$. In that version the exterior fibres need not be singletons.

For a cover $X=U_1\cup U_2$ by sieves, there is a second, particularly transparent argument in ordinary homotopy theory. Every simplex $x_0\to\cdots\to x_n$ in $NX$ lies in a member of the cover containing $x_n$, because a sieve contains every predecessor of its objects. Therefore
\[
 NX=NU_1\cup_{N(U_1\cap U_2)}NU_2
\]
is an equality of simplicial sets. The two inclusions are monomorphisms, so its realization is a pushout along subcomplex inclusions and computes the homotopy pushout. For cosieves one uses the first vertex $x_0$ instead. This calculation explains why categorical open and closed covers are especially well adapted to excision.

As a worked example, let $X$ be the poset of nonempty proper faces of a tetrahedron. Let $U_1$ consist of the faces of one facet, and let $U_2$ consist of the faces of the other three facets. Both are downward closed, and their union is $X$. Their intersection is the face poset of the boundary of the first facet. The order complexes are barycentric subdivisions: $BU_1$ is one triangle, $BU_2$ is the disk made of the other three triangles, and $B(U_1\cap U_2)$ is their common boundary circle. The integral therefore has homotopy type
\[
 D^2\mathop{\cup}^{h}_{S^1}D^2\simeq S^2,
\]
and the canonical comparison identifies this sphere with $BX$, the subdivided tetrahedral boundary. The comparison is determined by the actual inclusions of face posets; it is not merely an abstract equivalence between two spaces known to be spheres.
\end{dcomment}


Il suffit, en vertu du corollaire 2, de prouver que le carré est cocartésien dans Cat, puisque $i/t$ est une immersion ouverte ($i$ l'étant), et $X/t\longrightarrow X$ est un isomorphisme local, donc $W$-lisse. On voit que $(X/t)\backslash(U/t)\simeq e$ (catégorie ponctuelle), d'ailleurs l'unique élément de $(X/t)\backslash(U/t)$ est l'élément final $t_0 : t\xrightarrow{\mathrm{id}_t} t$ de $X/t$. Ainsi $H_{t_0}\simeq e_{U/t}$ (préfaisceau final), et $p_{!}(H_{t_0})\simeq H'_t \overset{\text{déf}}{=} (u\longmapsto \mathrm{Hom}_X(u,t))$, ce qui caractérise la somme amalgamée dans Cat
\dots

Le corollaire 4 n'est autre que le corollaire non démontré de la proposition 7 (p. 85).

Voici une variante du corollaire 4.

\noindent\textbf{Corollaire 5.} (Moyennant W(7 bis) (\footnote[105]{forme forte.}).) \textit{Soit $i : U \hto X$ une immersion ouverte, soit $t_0\in\mathrm{Ob}\,T$ un point \textup{fermé} de $T=X\backslash U$, tel que $\{t_0\}$}
\msp{page 151}
\textit{soit isomorphe à la catégorie ponctuelle, i.e. $\mathrm{Hom}_X(t_0,t_0)=\{\mathrm{id}\}$. Supposons que pour tout $t\in T\backslash\{t_0\}$, $U/t$ soit $W$-asphérique. Pour que l'on ait $i\in W$, il faut et il suffit que $i$ soit asphérique, i.e. que $U/t_0$ soit également $W$-asphérique.}

Il faut prouver que $i\in W$ implique que $U/t_0$ est $W$-asphérique. Considérons l'inclusion
\[
i_0 : U\longrightarrow U' \overset{\text{déf}}{=} X\backslash\{t_0\},
\]
l'hypothèse sur les $U/t$, pour $t\in T\backslash\{t_0\}$, signifie que $i_0$ est $W$-asphérique. Donc $i_0$ est dans $W$, et comme le composé $i : U\xrightarrow{i_0} U'\xrightarrow{i'} X$ est dans $W$, cela implique que
% foot: Version 7 décembre 2002
\printedpage{104}
% head: A. Grothendieck : Dérivateurs, Ch. XII, Édition des Universités de Montpellier II et Paris VII
$i' : U'\longrightarrow X$ est dans $W$. Nous sommes maintenant dans la situation du corollaire 4, qui implique que $i'$ est $W$-asphérique. Donc $i=i'i_0$ est $W$-asphérique, comme composé de deux morphismes $W$-asphériques, q.e.d.

\dgrblock{dgrb472}{A single missing point and the role of the strong axiom}
\begin{dcomment}{exposition}{p151:comme composé de deux morphismes $W$-asphériques}{A single missing point and the role of the strong axiom}
Grothendieck uses the preceding criterion to finish a result deferred earlier in the chapter. When only one object is missing, the slice over that object records every possible attachment to the open subcategory. The square on manuscript page 150 is consequently both a strict pushout and a homotopy pushout. The remaining implication is a reflection statement, and this is where the strong form of W(7 bis) enters.

Here is the strict pushout assertion in detail. Suppose $U\subset X$ is a sieve, its complement consists of one object $t$, and $\operatorname{End}_X(t)=\{\mathrm{id}_t\}$. Let $F:U\to Z$ and $G:X/t\to Z$ agree on $U/t$ after projection to $U$. Define $H:X\to Z$ on $U$ to be $F$, and set $H(t)=G(\mathrm{id}_t)$. For an arrow $a:u\to t$, define $H(a)$ to be the image under $G$ of the canonical arrow from $(u,a)$ to $(t,\mathrm{id}_t)$ in $X/t$. The source of this image is $F(u)$ by the agreement assumption. If $b:u'\to u$ lies in $U$, the factorization of $(u',ab)\to(t,\mathrm{id}_t)$ through $(u,a)$ gives
\[
 H(ab)=H(a)F(b).
\]
There are no arrows from $t$ into $U$, since $U$ is a sieve, and the only endomorphism at $t$ is its identity. Thus these rules define a functor, and they determine it uniquely. This proves the pushout universal property of the slice square, without any appeal to the presheaf notation $H_{t_0}$ used in the manuscript.

The projection $X/t\to X$ is a discrete fibration, hence smooth. Corollary 2 proves that this square is homotopy cocartesian without an extra excision hypothesis. If the lower inclusion $U\to X$ is in $W$, however, deducing that the upper inclusion $U/t\to X/t$ is in $W$ is exactly a converse-excision step. Under strong W(7 bis), Proposition 13 supplies it. The target $X/t$ has a final object, so $U/t$ is then $W$-aspherical. Conversely, if $U/t$ is $W$-aspherical, the inclusion $U\to X$ is $W$-aspherical: its slices at objects of $U$ have final objects, and its only remaining slice is $U/t$. Theorem A gives the weak equivalence.

The reduction in Corollary 5 is now precise. Because $t_0$ is closed in $X\setminus U$ and $U$ is a sieve, $t_0$ has no arrow to any other object of $X$. Hence $U'=X\setminus\{t_0\}$ is again a sieve. The hypotheses make $U\to U'$ aspherical. If $U\to X$ is in $W$, two-out-of-three makes $U'\to X$ a weak equivalence; Corollary 4 makes it aspherical. Composition then makes $U\to X$ aspherical, which supplies the remaining slice condition at $t_0$. The argument is complete under its stated strong hypothesis. It should not be read as a general converse to Theorem A for ordinary weak homotopy equivalences.
\end{dcomment}


\noindent\textsc{Démonstration de th. 4.}

\noindent\circled{a} Critère de lissité de $\varphi_Q : \mathcal{X} = \intdiag{X_0}{X_1}{X_2} \longrightarrow Y$.

Comme $\{1\}$, $\{2\}$ sont des ouverts de $\Phi = \begin{tikzcd}[row sep=small,column sep=small,baseline=(\tikzcdmatrixname-2-1.base)]
 & 1 \arrow[dl] \\
 0 & \\
 & 2 \arrow[ul]
\end{tikzcd}$, $i_1 : X_1\longrightarrow\mathcal{X}$ et $i_2 : X_2\longrightarrow\mathcal{X}$ sont des immersions ouvertes. Comme $\varphi_Q\circ i_1 = f_1$,
\msp{page 152}
$\varphi_Q\circ i_2=f_2$, on voit déjà que
\[
f_1 \text{ et } f_2 \text{ lisses} \iff \varphi_Q \text{ lisse sur l'ouvert réunion de } i_1(X_1) \text{ et } i_2(X_2)\text{, i.e. sur } \mathcal{X}\backslash i_0(X_0).
\]
Il reste à exprimer la lissité de $\varphi_Q$ en les points $x_0$ de $i_0(X_0)$. Soient $x_1$, $x_2$, $y$ les images de $x_0$ dans $X_1$, $X_2$, $Y$ respectivement. Il faut exprimer que les fibres de
\[
(\varphi_Q)/x_0 : \mathcal{X}/x_0\longrightarrow Y/y \tag*{$(*)$}
\]
sont asphériques. Nous allons pour cela expliciter $\mathcal{X}/x_0$ et la flèche $(*)$ en termes du carré commutatif
\[
(**)\qquad\qquad
\begin{tikzcd}[row sep=large,column sep=large]
X_0/x_0 \arrow[r, "\alpha_1/x_0"] \arrow[d, "\alpha_2/x_0"'] & X_1/x_1 \arrow[d, "f_1/x_1"] \\
X_2/x_2 \arrow[r, "f_2/x_2"] & Y/y,
\end{tikzcd}
\]
induit par $Q$. Je dis

\noindent\textbf{Lemme 1.} \textit{Soit un carré commutatif}
\[
Q \quad = \quad
\begin{tikzcd}[row sep=large,column sep=large]
X_0 \arrow[r, "\alpha_1"] \arrow[d, "\alpha_2"'] & X_1 \arrow[d, "f_1"] \\
X_2 \arrow[r, "f_2"] & Y
\end{tikzcd}
\]
\textit{dans} Cat, \textit{$x_0\in\mathrm{Ob}\,X_0$, $x_1$, $x_2$, $y$ ses images dans $X_1$, $X_2$, $Y$ respectivement, $Q/x_0$ le carré $(**)$ ci-dessus, induit par $Q$ sur les localisées $X_0/x_0$ etc., $\mathcal{X} = \intdiag{X_0}{X_1}{X_2}$. On a alors un}
% foot: Version 7 décembre 2002
\printedpage{105}
% head: A. Grothendieck : Dérivateurs, Ch. XII, Édition des Universités de Montpellier II et Paris VII
\msp{page 153}
\textit{isomorphisme canonique}
\[
\mathcal{X}/x_0 \quad\simeq\quad \intdiag{X_0/x_0}{X_1/x_1}{X_2/x_2} ,
\]
\textit{et la flèche $(\varphi_Q)/x_0$, localisée de $\varphi_Q : \mathcal{X}\longrightarrow Y$ en $x_0$, s'identifie à la flèche $\varphi_{(Q/x_0)}$ :}
\[
\varphi_{(Q/x_0)} : \quad \intdiag{X_0/x_0}{X_1/x_1}{X_2/x_2} \qquad\longrightarrow\qquad Y/y.
\]

\noindent\textsc{Démonstration} du Lemme. On sait que
\[
\mathcal{X}/x_0\longrightarrow\mathcal{X}
\]
est Cat-fibrée (scindée) (à fibres discrètes), et comme
\[
\mathcal{X}\longrightarrow\Phi
\]
est aussi Cat-fibrée (scindée), il s'ensuit qu'il en est de même de la flèche composée
\[
\mathcal{X}/x_0\longrightarrow\Phi.
\]
Donc pour `calculer' $\mathcal{X}/x_0$, en tant que catégorie sur $\Phi$, il faut expliciter ses trois fibres, et les deux foncteurs de changement de base qui les relient. On trouve aussitôt que ces fibres sont canoniquement isomorphes à
\[
X_0/x_0, \qquad X_1/x_1, \qquad X_2/x_2,
\]
et on admettra (!) que les foncteurs changement de base
\msp{page 154}
\[
X_0/x_0\longrightarrow X_1/x_1, \qquad X_0/x_0\longrightarrow X_2/x_2
\]
(définis à isomorphisme unique près) ne sont autres que les localisés $\alpha_1/x_0$ et $\alpha_2/x_0$ de $\alpha_1 : X_0\longrightarrow X_1$ et $\alpha_2 : X_0\longrightarrow X_2$. Cela montre donc qu'on a un isomorphisme
\[
\mathcal{X}/x_0 \quad\simeq\quad \intdiag{X_0/x_0}{X_1/x_1}{X_2/x_2} ,
\]
% foot: Version 7 décembre 2002
\printedpage{106}
% head: A. Grothendieck : Dérivateurs, Ch. XII, Édition des Universités de Montpellier II et Paris VII
et l'on admettra aussi que la flèche $(\varphi_Q)/x_0 : \mathcal{X}/x_0\longrightarrow Y/y$ s'identifie à $\varphi_{(Q/x_0)}$.

Pour exprimer la $W$-lissité de $\varphi_Q$ en $x_0\in\mathrm{Ob}\,X_0\subseteq\mathrm{Ob}\,\mathcal{X}$, on est donc ramené à exprimer la $W$-asphéricité des fibres de
\[
\varphi_{(Q/x_0)} : \quad \intdiag{X_0/x_0}{X_1/x_1}{X_2/x_2} \qquad\longrightarrow\qquad Y/y.
\]
Pour expliciter ces fibres, on a besoin d'un autre lemme auxiliaire

\dgrblock{dgrb473}{The slice of the integral, with its arrows}
\begin{dcomment}{exposition}{p154:Pour expliciter ces fibres, on a besoin d'un autre lemme auxiliaire}{The slice of the integral, with its arrows}
Grothendieck is reducing smoothness to a calculation at one object. The reduction succeeds because the smooth integral remembers the arrows into the central stratum, rather than merely the three categories as disjoint pieces. We make this calculation explicit; it supplies both identifications which the proof asks us to admit.

Write $E$ for the smooth integral of $X_0\xrightarrow{\alpha_i}X_i$, $i=1,2$. Its objects are the tagged objects of the three categories. Arrows within a stratum are the original arrows, and the additional hom-sets are
\[
 \Hom_E((i,b),(0,a))=\Hom_{X_i}(b,\alpha_i a),\qquad i=1,2.
\]
There are no arrows in the reverse direction or between the two outer strata. Composition is composition in $X_i$, using $\alpha_i$ when a central arrow is involved. Thus the cartesian arrow over $i\to0$ has the direction
\[
 (i,\alpha_i a)\longrightarrow(0,a).
\]
The arrows printed in the opposite direction on manuscript page 155 are an edition slip, not a different convention for this integral. The functor $\varphi_Q:E\to Y$ sends an additional arrow represented by $h:b\to\alpha_i a$ to $f_i(h)$, whose target is $f_i\alpha_i(a)=f_0(a)$. In particular it sends the cartesian arrow just displayed to an identity. This is the property which the source notices it had left implicit.

Fix $x_0\in X_0$ and put $x_i=\alpha_i x_0$. An object of $E/x_0$ in the central stratum is precisely an arrow $g:a\to x_0$ in $X_0$. An object in stratum $i$ is an arrow $h:b\to x_i$ in $X_i$. A cross-stratum morphism from $h$ to $g$ is an arrow $k:b\to\alpha_i a$ satisfying
\[
       (\alpha_i g)k=h.
\]
This is exactly a morphism from $h$ to the image of $g$ in $X_i/x_i$. Within each stratum we have the ordinary slice morphisms. These descriptions identify objects, hom-sets, identities, and compositions, giving the asserted canonical isomorphism
\[
 E/x_0\cong\int_{\mathrm{li}}
       (X_0/x_0\longrightarrow X_1/x_1,\ X_0/x_0\longrightarrow X_2/x_2).
\]
Under it the functor to $Y/y$, $y=f_0x_0$, sends $g$ to $f_0g$ and $h$ to $f_i h$; on a cross-stratum arrow it is $f_i k$. This proves the compatibility with $\varphi_{Q/x_0}$ as well as the isomorphism of the source categories.

The same description proves Lemma 2 without a further coherence argument. The strict fibre over $y$ consists of objects mapping to $y$ and arrows mapping to $1_y$. A cross-stratum arrow belongs to this fibre exactly when its representing arrow in $X_i$ does. Consequently
\[
 E_y\cong\int_{\mathrm{li}}(X_{0y}\longrightarrow X_{1y},\ X_{0y}\longrightarrow X_{2y}).
\]
Strict fibres need not be full subcategories. For instance, the strict fibre of the identity functor of a one-object group category over its object is the terminal category, whereas the full subcategory on that object is the whole group category. Keeping the identity condition on arrows is essential in the present calculation.

For an arrow $u:y'\to y$, set
\[
 A_i(u)=(X_i/x_i)_u.
\]
Its objects are arrows $a\to x_i$ whose image is exactly $u$; its morphisms are commuting triangles whose remaining edge maps to $1_{y'}$. The preceding two identifications now give, with no choices,
\[
       (E/x_0)_u\cong\int_{\mathrm{li}}
          (A_0(u)\longrightarrow A_1(u),\ A_0(u)\longrightarrow A_2(u)).
\]
This is the precise local input to the smoothness criterion. For the later systematic treatment, the locality of smoothness under passage to slices is Proposition 4.10 of \citet{Maltsiniotis2005Asphericite}; the argument here retains the additional explicit description needed for this particular integral.
\end{dcomment}


\noindent\textbf{Lemme 2} (\footnote[106]{NB C'est le b) du th. 4.}). \textit{Soit un carré commutatif dans} Cat
\[
Q \quad = \quad
\left(
\begin{tikzcd}[row sep=large,column sep=large]
X_0 \arrow[r, "\alpha_1"] \arrow[d, "\alpha_2"'] & X_1 \arrow[d, "f_1"] \\
X_2 \arrow[r, "f_2"] & Y
\end{tikzcd}
\right),
\]
\textit{d'où}
\[
\varphi_Q : \quad \left( \intdiag{X_0}{X_1}{X_2} \right) \qquad\longrightarrow\qquad Y,
\]

\msp{page 155}
\textit{et soit $y\in\mathrm{Ob}\,Y$. Alors la fibre de $\varphi_Q$ en $y$}
\[
\mathcal{X}_y = \varphi_Q^{-1}(\ \underbrace{\{y\}}_{\text{catégorie ponctuelle}}\ )
\]
\textit{est donnée par (comme catégorie au dessus de $\Phi$)}
\[
\mathcal{X}_y \quad\simeq\quad \intdiag{X_{0\,y}}{X_{1\,y}}{X_{2\,y}} .
\]

\noindent\textsc{Démonstration.} Il est immédiat qu'on a
\[
\mathcal{X}_y\cap\underbrace{\mathcal{X}_s}_{X_s}\simeq(X_s)_y \qquad\quad \text{pour } s\in\mathrm{Ob}\,\Phi=\{0,1,2\},
\]
% foot: Version 7 décembre 2002
\printedpage{107}
% head: A. Grothendieck : Dérivateurs, Ch. XII, Édition des Universités de Montpellier II et Paris VII
et (le hasard fait bien les choses) les foncteurs changement de base
\[
X_0\longrightarrow X_1, \qquad X_0\longrightarrow X_2
\]
envoient la sous-catégorie (non pleine) $X_{0\,y}$ de $X_0$ dans les sous-catégories (non pleines) $X_{1\,y}$ et $X_{2\,y}$ de $X_1$ et $X_2$ respectivement.
\[
\begin{tikzcd}[row sep=large,column sep=large]
X_{s\,y} \arrow[r] \arrow[d] \arrow[dr, phantom, "\scriptstyle\text{cart.}"] & \mathcal{X}_y \arrow[r] \arrow[d] \arrow[dr, phantom, "\scriptstyle\text{cart.}"] & \{y\} \arrow[d] \\
X_s \arrow[r] \arrow[d] \arrow[dr, phantom, "\scriptstyle\text{cart.}"] & \mathcal{X} \arrow[r] \arrow[d] & Y \\
\{s\} \arrow[r] & \Phi &
\end{tikzcd}
\]
On a donc une sous-catégorie fibrée de $\mathcal{X}$, sous réserve de vérifier que les flèches strictement cartésiennes
\[
x_0\longrightarrow\alpha_1(x_0), \qquad x_0\longrightarrow\alpha_2(x_0)
\]
dans $\mathcal{X}$, pour $x_0$ dans $X_{0\,y}$, sont dans $\mathcal{X}_y$, i.e. sont au-dessus de la flèche identique de $y$. Or c'est là la propriété essentielle
\msp{page 156}
du foncteur $\varphi_Q : \mathcal{X}\longrightarrow Y$, qu'on a oublié d'expliciter dans 3.2. Ça a vraiment l'air de marcher (comme sur des roulettes!)

Appliquons le Lemme 2 au diagramme carré $Q/x_0$ et à $\varphi_{(Q/x_0)}$ (page 154), on voit donc que les fibres à examiner en les $u\in\mathrm{Ob}\,Y/y$ ($u : y'\longrightarrow y$) peuvent s'écrire comme les intégrales
\[
\mathcal{X}(x_0,u) \quad\overset{\text{déf}}{=}\quad \intdiag{(X_0/x_0)_u}{(X_1/x_1)_u}{(X_2/x_2)_u} .
\]
Utilisons maintenant le fait que $f_1$, $f_2$ sont $W$-lisses (exprimant que $\varphi_Q$ est $W$-lisse dans $\mathcal{X}\backslash X_0$), donc que $(X_1/x_1)_u$ et $(X_2/x_2)_u$ sont $W$-asphériques. Dans le carré $W$-cocartésien
\[
\begin{tikzcd}[row sep=large,column sep=small]
 & (X_1/x_1)_u \arrow[dr, "{i_1(x_0,u)}"] & \\
(X_0/x_0)_u \arrow[ur, "{\alpha_1(x_0,u)}"] \arrow[dr, "{\alpha_2(x_0,u)}"'] & & (\mathcal{X}/x_0)_u\simeq\mathcal{X}(x_0,u) \\
 & (X_2/x_2)_u \arrow[ur, "{i_2(x_0,u)}"'] &
\end{tikzcd}
\]
% foot: Version 7 décembre 2002

% ---- from Der-ChXIIb-012_p066-070.tex
\printedpage{108}
% head: A. Grothendieck : Dérivateurs, Ch. XII, Édition des Universités de Montpellier II et Paris VII
(\footnote[107]{attention, ce n'est pas un carré commutatif. Dire plutôt que $(\mathcal{X}/x_0)_u \simeq \int \begin{tikzcd}[ampersand replacement=\&, row sep=0.4em, column sep=0.4em, baseline=(\tikzcdmatrixname-2-1.base)] \& {} \\ {} \ar[ur]\ar[dr] \& \\ \& {} \end{tikzcd}$ est lisse sur $\Phi$ à fibres $W$-asphériques (si $f_0$, $f_1$, $f_2$ sont lisses), donc $W$-asphérique puisque $\Phi$ l'est.}), il faut exprimer que $(\mathcal{X}/x_0)_u$ est $W$-asphérique, sachant que $(X_1/x_1)_u$ et $(X_2/x_2)_u$ le sont. Donc cela équivaut à % Resolved from PDF leaf 31: the original diagram has exactly these two unlabelled arrows and no vertex labels.
\[
i_1(x_0,u) \in W \qquad (\text{ou } i_2(x_0,u) \in W),
\]
donc aussi (moyennant W(7 bis)) à
\msp{page 157}
\[
\alpha_1(x_0,u) \in W, \qquad \text{i.e. } (X_0/x_0)_u \ W\text{-asphérique.}
\]
Dire que c'est vrai pour tout $u$ dans $Y/y$, signifie que $f_0 : X_0 \longrightarrow Y$ est lisse en $u$ \texttt{[plutôt en $x_0$]}.

Donc on a prouvé a), sous la forme plus précise :

\dgrblock{dgrb474}{Repairing the noncommuting fibre square}
\begin{dcomment}{exposition}{p157:Donc on a prouvé a), sous la forme plus précise :}{Repairing the noncommuting fibre square}
The aim of the calculation on page 156 is to reconstruct smoothness at the central object from smoothness on the three pieces. The displayed diamond cannot itself be used as a commutative square: its two composites land in different outer strata of the integral. They are related to the inclusion of the central stratum by cartesian natural transformations. Footnote 107 correctly withdraws that use of the square.

Here is the direct argument, with its hypotheses visible. Suppose $f_i$ is $W$-smooth at $x_i$ for $i=0,1,2$. For every $u:y'\to y$, all three categories $A_i(u)$ from the preceding block are $W$-aspherical. Their smooth integral $L_u$ is a split fibration over $\Phi=(1\to0\leftarrow2)$, hence is $W$-smooth. For each $s\in\Phi$, smoothness identifies its strict fibre, up to a $W$-equivalence, with the comma category $s\backslash L_u$. These comma categories are therefore $W$-aspherical. The dual form of the relative Theorem A gives $L_u\to\Phi$ in $W$. Since $0$ is final in $\Phi$, $L_u$ is $W$-aspherical. The slice calculation identifies $L_u$ with $(E/x_0)_u$, proving smoothness of $\varphi_Q$ at $x_0$. No pushout reflection axiom has entered.

For the converse, suppose the two outer maps are smooth at $x_1,x_2$, and retain the \emph{strong} form of W(7 bis) stipulated in the source. Then $A_1(u)$ and $A_2(u)$ are $W$-aspherical. Proposition 13, applied to their integral and its closed mapping-cylinder replacement, gives
\[
 A_1(u)\longrightarrow L_u\text{ in }W
 \quad\Longleftrightarrow\quad
 A_0(u)\longrightarrow A_2(u)\text{ in }W.
\]
The left side is equivalent to asphericity of $L_u$, and the right side to asphericity of $A_0(u)$, by two-out-of-three. This proves the asserted equivalence. It applies Proposition 13 to the correctly constructed integral; it does not promote the original noncommuting diamond to a strict square. The reverse implication is a reflection statement, stronger than the preservation statement used in the direct proof. The counterexample to strong reflection discussed earlier in the commentary on this chapter explains why it cannot be transferred without its hypothesis to ordinary weak homotopy equivalences.

At an outer object $x_i$, every arrow with target $x_i$ already belongs to $X_i$, since that stratum is open. Hence $E/x_i=X_i/x_i$, and local smoothness is exactly that of $f_i$. At a central object the hypotheses on both outer objects remain necessary for the converse equivalence just established. The final sentence after Corollary 6 drops them, as the editorial annotation observes; the following finite example measures precisely what is lost.
\end{dcomment}


\textbf{Corollaire 6.} \textit{Soit $x \in \Ob \mathcal{X}$ (où $\mathcal{X} = \int \begin{tikzcd}[ampersand replacement=\&, row sep=tiny, column sep=small, baseline=(\tikzcdmatrixname-2-1.base)] \& X_1 \\ X_0 \ar[ur]\ar[dr] \& \\ \& X_2 \end{tikzcd}$.) On veut exprimer la $W$-lissité de}
\[
\varphi_Q : \mathcal{X} \longrightarrow Y
\]
\textit{en $x$.}

\begin{enumerate}[label={}, leftmargin=2.5em]
\item[\textit{a)}] \textit{Si $x \in \Ob X_i$ ($i \in \{1,2\}$), $\varphi_Q$ est $W$-lisse en $x$ si et seulement si $f_i : X_i \longrightarrow Y$ l'est.}
\item[\textit{b)}] \textit{Supposons} \textup{W(7 bis)} (\footnote[108]{Inutile pour la partie directe, en supposant \texttt{[phrase incomplète]}.})\textit{. Soit $x_0 \in \Ob X_0$, et supposons $f_1$, $f_2$ lisses en $x_1$, $x_2$. Alors $\varphi_Q$ est $W$-lisse en $x_0$ si et seulement si $f_0 : X_0 \longrightarrow Y$ l'est.}
\end{enumerate}

\textit{Donc} (moyennant W(7 bis)) \textit{dans les trois cas $x \in \Ob X_i$ ($i \in \{0,1,2\}$) on trouve que $\varphi_Q$ est $W$-lisse en $x$ si et seulement si $f_i : X_i \longrightarrow Y$ l'est} \texttt{[pas vraiment pour $i=0$]}.

\cercle{b} Quant à b) du théorème 4, c'est le lemme 2 ci-dessus. Donc tout est prouvé!

\msp{page 158}

\dgrblock{dgrb475}{A finite witness for the missing outer hypothesis}
\begin{dcomment}{development}{p157}{A finite witness for the missing outer hypothesis}
Grothendieck has isolated a useful local criterion, but the final formulation would make the middle map alone control the whole neighbourhood. We give a small witness showing why the two outer hypotheses must stay in the statement. We did not find this particular finite witness in the smoothness discussion of \citet{Maltsiniotis2005Asphericite} or the integration discussion of \citet[\S3.3]{Cisinski2006}; it is an elementary application of their criteria, not a claim of priority.

\textbf{Proposition (failure of the unconditional central criterion).}\ours
Let $W$ be a fundamental localizer satisfying W(5), so that every $W$-equivalence induces a bijection on connected components. There is a commutative square of finite posets for which $f_0$ is smooth and $f_2$ is smooth, but the smooth integral map $\varphi_Q$ is not $W$-smooth at an object of the central stratum.

\textit{Proof.} Let $Y=X_0=X_2=[1]=(0\to1)$, with $f_0=f_2=\alpha_2=1$. Let $X_1$ be the poset with objects $a,b,c$ and only nonidentity relations $a<c$ and $b<c$. Define
\[
 f_1(a)=f_1(b)=0,\qquad f_1(c)=1,
 \qquad \alpha_1(0)=a,\quad\alpha_1(1)=c.
\]
Both nonidentity arrows of $X_1$ map to $u:0\to1$. These formulas define functors and give $f_1\alpha_1=1_Y$, so the required square commutes. Identity functors are smooth; this proves the assertions about $f_0$ and $f_2$.

Take $x_0=1$, $x_1=c$, $x_2=1$ and the arrow $u:0\to1$ in $Y$. The three strict lift categories are
\[
 A_0(u)=e,\qquad A_1(u)=\{a\to c,\ b\to c\}_{\mathrm{disc}},
 \qquad A_2(u)=e.
\]
The first map chooses $a\to c$, and the second is the unique map between terminal categories. Their smooth integral has four objects: a central object $z$, outer objects $a_*,b_*$ in stratum 1, and $d_*$ in stratum 2. Its only nonidentity arrows are
\[
                 a_*\longrightarrow z\longleftarrow d_*.
\]
The object $b_*$ is isolated. Thus its connected-component set has two elements. By W(5) it is not $W$-aspherical. The slice identification gives this category as $(E/x_0)_u$, which violates the defining local smoothness test.\hfill$\square$

The failure occurs in a missing lift relation at the first outer stratum. It involves no question about strictification and no exotic homotopy type. In Corollary 6(b), the assumption that $f_1$ is smooth at $x_1$ excludes exactly this additional component.
\end{dcomment}


\subsection*{3.5\quad Applications aux sommes amalgamées dans une catégorie $X^{\wedge}$ (avec $X \in \Ob \Cat$).}

D'abord une caractérisation des sommes amalgamées impliquant une immersion ouverte :

\textbf{Proposition 15.} \textit{Soit}
% foot: Version 7 décembre 2002
\printedpage{109}
% head: A. Grothendieck : Dérivateurs, Ch. XII, Édition des Universités de Montpellier II et Paris VII
\[
Q : \qquad
\begin{tikzcd}
U \ar[r, "i"] \ar[d, "p"'] & X \ar[d, "q"] \\
U' \ar[r, "i'"] & X'
\end{tikzcd}
\]
\textit{un carré commutatif dans} $\Cat$\textit{. Pour que ce carré soit cocartésien dans} $\Cat$\textit{, avec $i$ (donc $i'$) une immersion ouverte, il faut et il suffit que les conditions suivantes soient satisfaites :}

\begin{enumerate}[label={}, leftmargin=2.5em]
\item[\textit{a)}] \textit{$Q$ est cartésien.}
\item[\textit{b)}] \textit{$i'$ est une immersion ouverte (donc $i$ aussi).}
\item[\textit{c)}] \textit{Si $T' = X' \backslash U'$, et $T = X \times_{X'} T' = X \backslash U$, $T \longrightarrow T'$ est un isomorphisme.}
\item[\textit{d)}] \textit{Pour tout $t \in T$, si $t'$ est son image dans $T'$, le foncteur induit par $p$}
\[
p/t : \underbrace{U/t}_{= V_t} \longrightarrow \underbrace{U'/t'}_{= V'_{t'}}
\]
\textit{satisfait à $(p/t)_!(e_{U/t}) \simeq e_{U'/t'}$, ou ce qui revient au même, les `fibres colocales' de $p/t$, $v' \backslash V_t$, avec $v'$ dans $V'_{t'}$, sont} $0$\textit{-connexes.}
\end{enumerate}

\msp{page 159}
Démonstration immédiate - ça devrait plutôt remonter au §1.

\dgrblock{dgrb476}{Gluing the morphisms across an open immersion}
\begin{dcomment}{exposition}{p158}{Gluing the morphisms across an open immersion}
Grothendieck now changes scale: the homotopical assertion will rest on an ordinary categorical pushout. An open immersion $U\hookrightarrow X$ is a full sieve. Its complementary full subcategory $T$ is a cosieve; there are no arrows from $T$ to $U$. Thus $X$ is specified by $U$, $T$, and the bifunctor of crossing arrows
\[
       H:U^{\mathrm{op}}\times T\longrightarrow\mathrm{Set},
       \qquad H(u,t)=\Hom_X(u,t).
\]
Precomposition and postcomposition encode every composition involving a crossing arrow. This description explains the four conditions in Proposition 15: the first three identify the two pieces, while the fourth identifies the arrows joining them.

Assume first that $Q$ is cartesian, that $U'\hookrightarrow X'$ is open, and that the induced map $T\to T'$ is an isomorphism. Identify $T$ with $T'$ by this specified map. For $t\in T$, let $H_t$ and $H'_t$ be the presheaves $\Hom_X(-,t)$ and $\Hom_{X'}(-,t)$ on $U$ and $U'$. The square supplies a natural comparison
\[
                  p_!H_t\longrightarrow H'_t.
\]
The universal category obtained by pushing $U\hookrightarrow X$ along $p$ has pieces $U'$ and $T$ and crossing-arrow presheaf $p_!H_t$. Indeed, a functor from this category to $Z$ is a pair of functors on $U'$ and $T$, together with compatible images of crossing arrows; by the defining adjunction of $p_!$, these data are exactly a functor $X\to Z$ and a functor $U'\to Z$ agreeing on $U$. Composition with arrows of $T$ is preserved because $H_t$ varies functorially with $t$. This proves the universal property, not merely an objectwise description. Consequently $Q$ is a pushout precisely when all the displayed comparisons are isomorphisms.

We can test these comparisons by explicit factorization categories. An element of $(p_!H_t)(u')$ is represented by a pair
\[
       b:u'\longrightarrow p(u),\qquad a:u\longrightarrow t,
\]
modulo the relation generated by morphisms $h:u\to v$ with $a=a_vh$ and $p(h)b=b_v$. Its image is the composite $q(a)b:u'\to t$ in $X'$. Fix such a composite $a':u'\to t$. The category of its representatives and these morphisms is exactly
\[
        (u',a')\backslash(U/t),
\]
where $U/t\to U'/t$ is induced by $p$. Thus the fibre of the comparison over $a'$ is the set of connected components of this category. The comparison is bijective for every $u'$ if and only if every such category is nonempty and connected. This proves condition (d), including the equivalence with $(p/t)_!e=e$.

Conversely, the explicit universal category just constructed has $U'$ as a full sieve, has complement exactly $T$, and its inverse image of $U'$ is exactly $U$. Hence a pushout along $U\hookrightarrow X$ necessarily satisfies (a)--(c); its crossing-arrow comparison is an isomorphism, so it also satisfies (d). This completes both directions of Proposition 15. The argument is entirely about ordinary categories and ordinary left Kan extension.

There is a small example showing why the last condition cannot be omitted. Take $U=U'=e$, $X$ the discrete category on $u,t$, and $X'=(u\to t)$, with $q$ the identity on objects and $p=1$. The square is cartesian; both open subcategories are $\{u\}$, and the complements are the same terminal category. But $U/t$ is empty whereas $U'/t=e$. The missing crossing arrow makes the factorization category empty, so (d) fails. The pushout along the identity $p$ is $X$, not $X'$. The condition about connected factorization categories is exactly the information which the object sets and the two strata do not contain.
\end{dcomment}


\textbf{Corollaire 1.} \textit{Soit $Q$ un carré} cartésien (cf. page 158), \textit{avec $i'$ une} immersion ouverte \textit{(donc $i$ aussi), et $q$} $W$-lisse \textit{(donc $p$ aussi). On suppose que $W$ satisfait l'axiome} W(5) \textit{de connexité} (\footnote[109]{Donc il suffit $q$ $W_0$-lisse, où $W_0 = \{f \in \Fl \Cat \mid \pi_0(f) \text{ bijectif}\}$.})\textit{. Soit $T' = X' \backslash U'$, $T = X \backslash U$, supposons que $T' \longrightarrow T$ soit un isomorphisme. Alors $Q$ est cocartésien.}

Il suffit de prouver la condition d) de prop. 15. Considérons le carré
\[
\begin{tikzcd}
U/t \ar[r, hook, "i/t"] \ar[d, "p/t"'] & X/t \ar[d, "q/t"] \\
U'/t' \ar[r, hook, "i'/t'"] & X'/t',
\end{tikzcd}
\tag*{$(*)$}
\]
je dis qu'il est \textit{cartésien.} Cela résulte de la seule hypothèse que $Q$ soit cartésien, sans hypothèse de lissité sur $q$ ni que $i'$ soit une immersion ouverte, et se voit p. ex. en regardant le diagramme cubique :
% foot: Version 7 décembre 2002
\printedpage{110}
% head: A. Grothendieck : Dérivateurs, Ch. XII, Édition des Universités de Montpellier II et Paris VII
\[
\begin{tikzcd}[row sep=small, column sep=small]
& U/t \ar[rr] \ar[dl] \ar[dd] & & X/t \ar[dl] \ar[dd] \\
U'/t' \ar[rr, crossing over] \ar[dd] & & X'/t' & \\
& U \ar[rr] \ar[dl] & & X \ar[dl] \\
U' \ar[rr] & & X' \ar[from=uu, crossing over] &
\end{tikzcd}
\]
(\footnote[110]{NB Ici $t$ est un élément quelconque de $\Ob X$, $t'$ son image dans $X'$.}). Dans ce diagramme, la face antérieure et postérieure du cube est cartésienne, par définition de $U/t$ et de $U'/t'$
\msp{page 160}
respectivement, et la base l'est par hypothèse. Il en résulte que la face supérieure du cube est également un carré cartésien.

Ceci vu, commme $q$ est lisse, il en est de même de $q/t$ dans le carré cartésien $(*)$, donc aussi de $p/t$ par changement de base. D'autre part, la lissité de $q$ implique que $q/t$ est à fibres $W$-asphériques, donc par changement de base il en est de même de $p/t$ (qui est donc une $W$-équivalence universelle, i.e. une $W$-fibration triviale). De plus, à cause de la lissité de $p/t$, pour $v'$ dans $U'/t'$, l'inclusion % [sic] commme -> comme, repeated letter
\[
(U/t)_{v'} \longrightarrow v' \backslash (U/t)
\]
est dans $W$, donc (par W(5)) définit une bijection sur les $\pi_0$. Comme $(U/t)_{v'}$ est $W$-asphérique par ce qui précède, il s'ensuit (par W(5)) que son $\pi_0$ est de cardinal 1, donc il en est de même de $v' \backslash (U/t)$, i.e. celui-ci est 0-connexe. Or c'est là la condition d) de la proposition 15, qu'il fallait établir.

\msp{page 161}
\textbf{Corollaire 2.} \textit{Soit $Q$ un carré cartésien} (p. 158), \textit{avec $i'$ (donc aussi $i$) une immersion ouverte, et $q$ (donc aussi $p$) un} isomorphisme local. \textit{Soit $T = X \backslash U$, $T' = X' \backslash U'$ (donc $T \simeq X \times_{X'} T'$), et supposons $T' \longrightarrow T$ isomorphisme. Alors $Q$ est cocartésien.}

C'est un cas particulier du corollaire 1, mais sans référence explicite à un localiseur $W$. On peut pendre $W = W_{\infty}$, ou $W = W_0$, pour se ramener au corollaire 1. Mais il est plus raisonnable de dire que dans le carré cartésien $(*)$, l'hypothèse que $q$ soit un isomorphisme local implique que $q/t$ est un isomorphisme, donc $p/t$ l'est aussi. Il en résulte bien que les $v' \backslash (U/t)$ sont 0-connexes, cette catégorie étant isomorphe à $v' \backslash (U'/t')$, donc a un objet initial. % [sic] pendre -> prendre, missing letter

\dgrblock{dgrb477}{From smoothness to the connected factorization test}
\begin{dcomment}{exposition}{p161:C'est un cas particulier du corollaire 1, mais sans référence explicite à un localiseur $W$.}{From smoothness to the connected factorization test}
The two corollaries of Proposition 15 replace its explicit factorization condition by familiar geometric hypotheses. Their force is that a condition on the ambient map $q$ controls every crossing arrow simultaneously.

For a cartesian square $Q$ and an object $t\in X$, the square of slices displayed on page 159 is cartesian. An object of its pullback is an arrow $x\to t$ in $X$ together with a chosen lift of its image-source to $U'$; the cartesian property of $Q$ identifies that choice with an object of $U$ lying above $x$. The same description of morphisms gives the isomorphism with $U/t$. No smoothness is needed for this step.

If $q$ is $W$-smooth, then $q/t:X/t\to X'/t'$ is $W$-smooth and has $W$-aspherical strict fibres. The first assertion is locality of smoothness on slices; the second is its defining lift-category test at $t$. Both statements are inherited by the pullback $p/t$. For $v'\in U'/t'$, smoothness gives a $W$-equivalence
\[
                (U/t)_{v'}\longrightarrow v'\backslash(U/t).
\]
Its source is $W$-aspherical, so its target is too. W(5) converts this into the required statement that the target has exactly one connected component. Proposition 15 then proves that $Q$ is cocartesian. Only the connectedness part of smoothness is used: as footnote 109 says, $W_0$-smoothness is sufficient.

For a local isomorphism $q$, the entire slice map $q/t$ is an isomorphism, and so is $p/t$. In that case every factorization category $v'\backslash(U/t)$ has an initial object. This gives the second corollary without any localizer. The later formulation that local isomorphisms are smooth is \citet[Proposition 4.9]{Maltsiniotis2005Asphericite}, but the direct slice argument also explains the stronger ordinary conclusion needed here.

The arrows on complements in the statements of both corollaries are printed $T'\to T$. The map induced by the square is $T\to T'$, as in Proposition 15(c) and in the application on page 166. The condition is that this induced map be an isomorphism. Its inverse then exists, but no inverse map is part of the initial data. We retain the printed direction in the text and record the correction in the register.
\end{dcomment}


Soit maintenant $X$ un objet de Cat, et considérons un diagramme cocartésien dans $X^{\wedge}$ :
\[
\begin{tikzcd}
G \ar[r, "i"] \ar[d, "p"'] & F \ar[d, "q"] \\
G' \ar[r, "i'"] & F',
\end{tikzcd}
\tag*{$(*)$}
\]
% foot: Version 7 décembre 2002
\printedpage{111}
% head: A. Grothendieck : Dérivateurs, Ch. XII, Édition des Universités de Montpellier II et Paris VII
\msp{page 162}
où $i$ est un monomorphisme. On sait que le caractère cocartésien du carré se vérifie `argument par argument' par le caractère cocartésien de
\[
\begin{tikzcd}
G(x) \ar[r, "i(x)"] \ar[d, "p(x)"'] & F(x) \ar[d, "q(x)"] \\
G'(x) \ar[r, "i'(x)"] & F'(x),
\end{tikzcd}
\tag*{$(*)$}
\]
dans Ens, $\forall\, x$ dans $X$. Or un tel carré cocartésien dans Ens, avec $i(x)$ \texttt{[un]} monomorphisme, est aussi cartésien. Cela implique que le carré $Q$ dans $X^{\wedge}$ est également cartésien.

Considérons, d'autre part, le foncteur canonique
\[
\varphi_X : F \longmapsto X/F : X^{\wedge} \longrightarrow \mathrm{Cat}.
\]
Il vaut mieux, d'ailleurs, le regarder comme un foncteur
\[
\varphi_X : X^{\wedge} \longrightarrow \mathrm{Cat}/X,
\]
et je rappelle, pour mémoire, que ce foncteur est pleinement fidèle, et que son image essentielle est formée des catégories $X'$ sur $X$ telles que
\[
X' \longrightarrow X
\]
\msp{page 163}
soit un isomorphisme local, ou ce qui revient au même, que $X'$ soit \textit{fibrée} sur $X$ à fibres discrètes :
\[
\varphi_X : X^{\wedge} \xrightarrow{\ \sim\ } \underbrace{\Fibdisc\, X}_{\substack{\text{catégories fibrées}\\ \text{à fibres discrètes}}} \ .
\]
Notons que toute flèche $X' \longrightarrow X''$ dans $\Fibdisc\, X$ est un isomorphisme local (tautologique sur $\begin{tikzcd}[ampersand replacement=\&, row sep=small, column sep=tiny, baseline=(\tikzcdmatrixname-1-1.base)] X'/x' \ar[rr] \ar[dr, "\text{iso.}"'] \& \& X''/x'' \ar[dl, "\text{iso.}"] \\ \& X/x \& \end{tikzcd}$), ce qui montre que si on a un diagramme
\[
\begin{tikzcd}[row sep=small, column sep=small]
X_1 \ar[dr] & & X_2 \ar[dl] \\
& X_0 &
\end{tikzcd}
\]
dans $\Fibdisc\, X$, le produit fibré $X_3 = X_1 \times_{X_0} X_2$ dans Cat est tel que $X_3 \longrightarrow X_2$ est un isomorphisme local, donc aussi $X_3 \longrightarrow X$. C'est-à-dire $\Fibdisc(X)$ est stable par produits fibrés. Ceci implique que les produits fibrés existent dans $\Fibdisc\, X$, et que ce sont les produits fibrés dans Cat, ou dans Cat$/X$.
% foot: Version 7 décembre 2002

\printedpage{112}
% head: A. Grothendieck : Dérivateurs, Ch. XII, Édition des Universités de Montpellier II et Paris VII
\msp{page 164}
De même pour les produits finis dans $\Fibdisc\, X$, ce sont les produits finis dans Cat$/X$. Bref

\textbf{Lemme 1.} \textit{La sous-catégorie pleine} $\Fibdisc\, X$ \textit{de} $\Cat/X$ \textit{est stable par} $\varprojlim$ \textit{finies, donc les} $\varprojlim$ \textit{existent dans} $\Fibdisc\, X$ \textit{et le foncteur d'inclusion}
\[
\Fibdisc\, X \longrightarrow \Cat/X
\]
\textit{y commute, et en particulier transforme carrés cartésiens de} $\Fibdisc\, X$ \textit{en carrés cartésiens de} $\Cat/X$\textit{, ou ce qui revient au même, de} $\Cat$\textit{. De plus, en vertu de l'équivalence}
\[
\varphi_X : X^{\wedge} \xrightarrow{\ \sim\ } \Fibdisc\, X \hookrightarrow \Cat/X,
\]
\textit{le foncteur $\varphi_X$ de $X^{\wedge}$ dans} $\Cat/X$ \textit{est exacte à gauche, en particulier il transforme carrés cartésiens de $X^{\wedge}$ en carrés cartésiens de} $\Cat/X$\textit{, ou encore de} $\Cat$\textit{.} % [sic] exacte -> exact, agreement with "le foncteur"

\msp{page 165}
Notons aussi qu'étant exact à gauche, le foncteur $\varphi_X$ de $X^{\wedge}$ dans Cat$/X$ transforme monomorphismes en monomorphismes. D'autre part :

\dgrblock{dgrb478}{The exactness of the category of elements}
\begin{dcomment}{later}{p164}{The exactness of the category of elements}
Grothendieck's finite-limit lemma is the first part of a stronger known statement. The relative category-of-elements functor
\[
             j_X:\widehat X\longrightarrow\Cat/X,
             \qquad F\longmapsto (X/F\to X),
\]
preserves all small limits and colimits. This is Proposition 3.2.12 of \citet{Cisinski2006}, in the published version, page 141. The absolute functor $i_X:F\mapsto X/F$ preserves colimits and pullbacks, and has the right adjoint stated in Corollary 3.2.13, page 142. We spell out the mechanism because it will also explain the pushout in Theorem 5.

An object of $X/F$ is $(x,a)$ with $a\in F(x)$. An arrow $(x,a)\to(y,b)$ is an arrow $u:x\to y$ such that $F(u)b=a$. Thus an $n$-simplex of its nerve is uniquely determined by a chain in $X$ and the element at its last vertex:
\[
 N(X/F)_n\cong
       \coprod_{x_0\to x_1\to\cdots\to x_n} F(x_n).
\]
All earlier elements are recovered by restriction. The final face uses the last restriction map of $F$; the other face and degeneracy maps change the underlying chain. The displayed bijections therefore define the nerve functor, not just its degreewise cardinalities.

Colimits of presheaves are evaluated pointwise, and coproducts commute with colimits. It follows that $N(X/F)$ preserves every small colimit in $F$. A cocone of nerves whose colimit is again a nerve is also a colimit cocone in categories: apply the fully faithful nerve to maps into an arbitrary target category. This proves the colimit assertion for $i_X$, and the cocones retain their maps to $X$, proving it for $j_X$ as well. For limits over $X$, work over the fixed simplicial set $NX$. The fibre in degree $n$ above a specified chain is $F(x_n)$, so limits are again calculated pointwise. Since the nerve creates limits, this proves the limit assertion for $j_X$.

For completeness, the absolute right adjoint is
\[
       R_X(C)(x)=\operatorname{Fun}(X/x,C),
\]
where the right side is the set of functors. For $u:x\to y$, restriction is precomposition with $X/x\to X/y$. A functor $H:X/F\to C$ assigns to $a\in F(x)$ the composite
\[
 X/x\longrightarrow X/F\xrightarrow{H}C,
 \qquad (z\xrightarrow{v}x)\longmapsto(z,F(v)a).
\]
These assignments form a map $F\to R_X(C)$. Conversely, given such a map, evaluate its functor at $1_x$ to obtain the image of $(x,a)$. An arrow $u:(x,a)\to(y,b)$ is sent to the arrow in the assigned functor on $X/y$ from $u$ to $1_y$. Naturality gives the correct source and ensures preservation of identities and compositions. The two constructions are inverse. This proves the adjunction directly.

The relative target matters for products: $X/(F\times G)$ is the fibre product $(X/F)\times_X(X/G)$. It is generally not the ordinary product of the two total categories. The absolute functor sends the terminal presheaf to $X$, not to the terminal category. The manuscript correctly formulates finite exactness in $\Cat/X$.
\end{dcomment}


\textbf{Lemme 2.} \textit{Les monomorphismes qui sont des isomorphismes locaux sont exactement les immersions ouvertes. Donc dans} $\Fibdisc\, X$\textit{, les monomorphismes sont les immersions ouvertes, et le foncteur $\varphi_X : F \longmapsto X/F$ transforme monomorphismes en immersions ouvertes.}

Revenons alors au diagramme cocartésien $(*)$ (p. 161), et au diagramme correspondant dans Cat (et même dans $\Fibdisc\, X$)
\[
\begin{tikzcd}
X/G \ar[r, hook, "X/i"] \ar[d, "X/p"'] & X/F \ar[d, "X/q"] & T \ar[l, hook', dotted] \ar[d, dotted] \\
X/G' \ar[r, hook, "X/i'"] & X/F' & T'. \ar[l, hook', dotted]
\end{tikzcd}
\]
Comme $(*)$ est cartésien, le carré précédent est cartésien, par ce qui précède. Je veux prouver qu'il est également cocartésien
\msp{page 166}
(dans Cat$/X$, ou dans Cat, cela revient au même), et pour ceci j'applique le corollaire 2 de la proposition 15 (p. 161), compte tenu que $X/q$ est un isomorphisme local, et $X/i'$ (étant un isomorphisme local \texttt{[et un]} monomorphisme) une immersion ouverte. Il suffit de prouver que, $T'$ et $T$ étant les sous-catégories fermées complémentaires des ouverts $X/G'$ de $X/F'$ et de $X/G$ de $X/F$, que $T \longrightarrow T'$ est \texttt{[un]} isomorphisme.

{\small \texttt{[Maintenant, $F$ et $G$, et $F'$ et $G'$, mais pas $p$ et $q$, échangent leur rôles.]}} % [sic] leur rôles -> leurs rôles, agreement

La vérification est immédiate : $T$ est la catégorie des flèches $x \longrightarrow G$ dans $X^{\wedge}$ (de source $x$ dans $X$) qui ne se factorisent pas par $G$ \texttt{[plutôt par $F$]}, $T'$ celle des flèches $x \longrightarrow G'$
% foot: Version 7 décembre 2002

% ---- from Der-ChXIIb-012_p071-072.tex
\printedpage{113}
% head: A. Grothendieck : Dérivateurs, Ch. XII, Édition des Universités de Montpellier II et Paris VII
qui ne se factorisent pas par $F'$, le foncteur $T \longrightarrow T'$ étant le foncteur composition de $x \longrightarrow G$ avec $G \xrightarrow{q} G'$
\[
\begin{tikzcd}[row sep=large]
F \arrow[r, hook] \arrow[d] & G \arrow[d] \\
F' \arrow[r, hook] & G'
\end{tikzcd}
\]
(NB Le composé ne se factorise pas par $F'$, car autrement, à cause du caractère cartésien de $Q$, $x \longrightarrow G$ se factoriserait par $F$.) Mais pour $x$ donné, la donnée de $x \longrightarrow G$ ne se factorisant pas par $F$ est celle d’un élément \msp{page 167} de $G(x)\backslash F(x)$, et à cause du caractère cocartésien de $Q$,
\[
(G(x)\backslash F(x)) \longrightarrow (G'(x)\backslash F'(x))
\]
est bijectif, i.e. ‘il revient au même’ de se donner $x \longrightarrow G$ ne se factorisant pas par $F$, ou $x \longrightarrow G'$ ne se factorisant pas par $F'$.

\dgrblock{dgrb479}{Why the complement calculation is enough}
\begin{dcomment}{exposition}{p167:est bijectif, i.e. ‘il revient au même’ de se donner}{Why the complement calculation is enough}
Grothendieck wants to transport a pushout of presheaves to a pushout of categories, where the homotopical criterion is already available. We use throughout the notation of the final theorem, $F\hookrightarrow G$ and $F'\hookrightarrow G'$. The source changes the letters immediately before this calculation; the editorial note marks that change, and no mathematical reversal is intended.

At each $x\in X$, the pushout of sets has the concrete description
\[
 G'(x)\cong F'(x)\amalg\bigl(G(x)\setminus F(x)\bigr).
\]
The map $q_x$ sends an element of $F(x)$ through $p_x$ and leaves the complementary elements distinct. Hence $F'(x)\to G'(x)$ is injective and its inverse image in $G(x)$ is exactly $F(x)$. This proves simultaneously that $i'$ is a monomorphism and that the square of presheaves is cartesian. It also gives a bijection between the two objectwise complements.

A monomorphism $F\hookrightarrow G$ induces a full sieve $X/F\hookrightarrow X/G$: if $(y,b)$ belongs to $X/F$ and $u:(x,a)\to(y,b)$ is an arrow of $X/G$, then $a=G(u)b$ belongs to $F(x)$. Conversely, a monomorphism which is a local isomorphism is a full sieve. To see this, lift uniquely every arrow into an object in its image using the isomorphism on slices; injectivity on objects and arrows then identifies the lift with the given arrow. This is Lemma 2 with its categorical content made explicit.

It remains to check that the bijection of complementary objects respects all arrows in both directions. Let $(x,a)$ and $(y,b)$ be outside $X/F$. An arrow between their images in $X/G'$ has an underlying arrow $u:x\to y$ satisfying
\[
       G'(u)q_y(b)=q_x(a).
\]
By naturality its left side is $q_xG(u)b$. The element $G(u)b$ cannot belong to $F(x)$, since then its image would belong to $F'(x)$, whereas $q_x(a)$ does not. Both $G(u)b$ and $a$ are therefore in the complement where $q_x$ is injective. They are equal. This gives precisely one arrow in the original complement, with the same underlying $u$. Identities and compositions are preserved because both complement functors lie over $X$. The functor $T\to T'$ is thus an isomorphism of categories, not merely a bijection on objects. Proposition 15, Corollary 2, now proves the categorical pushout.

The stronger exactness result in the previous block gives an alternative: $i_X$ preserves this pushout even without the monomorphism hypothesis. The monomorphism is needed for the homotopical conclusion, and for the simultaneous pullback assertion. Indeed, the nerve formula sends the square to a simplicial pushout along a monomorphism. Such a pushout computes the homotopy pushout of the nerves. Together with the Grothendieck-construction comparison, this proves that its canonical smooth-integral map is in $W_\infty$, and hence in every fundamental localizer, by \citet[Corollary 4.2.19]{Cisinski2006}. Thus the W(7 bis) proviso in the homotopy-pushout clause of Theorem 5 can be removed. It must still be retained for the \emph{reflection} implications in the following corollary.
\end{dcomment}


Ainsi on a prouvé le

\textbf{Théorème 5.} \textit{Soit $X$ dans} Cat\textit{, et}
\[
Q : \qquad
\begin{tikzcd}[row sep=large, baseline=(current bounding box.center)]
F \arrow[r, hook, "i"] \arrow[d, "p"'] & G \arrow[d, "q"] \\
F' \arrow[r, hook, "i'"] & G'
\end{tikzcd}
\]
\textit{un carré cocartésien dans $X^\wedge$, avec $i$} \texttt{[un]} \textit{monomorphisme. Alors $i'$ est} \texttt{[un]} \textit{monomorphisme, le carré $Q$ est aussi cartésien, et le carré correspondant dans} Cat$/X$ \textit{ou dans} Cat
\[
\begin{tikzcd}[row sep=large]
X/F \arrow[r, hook, "X/i"] \arrow[d, "X/p"'] & X/G \arrow[d, "X/q"] \\
X/F' \arrow[r, hook, "X/i'"] & X/G'
\end{tikzcd}
\]
\msp{page 168} \textit{est à la fois cartésien et cocartésien. De plus, si $W$ est un localiseur fondamental dans} Cat\textit{, satisfaisant} W(7 bis) (en plus de W(1,2,3))\textit{, alors le carré $X/Q$ dans} Cat \textit{est $W$-cocartésien.}

\printedpage{114}
% head: A. Grothendieck : Dérivateurs, Ch. XII, Édition des Universités de Montpellier II et Paris VII
Ce dernier énoncé résulte de ce qui précédait et du corollaire 2 du théorème 4, compte tenu que $X/q$ est un isomorphisme local, donc est $W$-lisse.

\dgrblock{dgrb480}{An attachment over the walking arrow}
\begin{dcomment}{examples}{p168:Ce dernier énoncé résulte de ce qui précédait}{An attachment over the walking arrow}
Theorem 5 can be seen without any large category of simplices. Let $X=[1]=(0\to1)$. A presheaf on $X$ is a map of sets $F(1)\to F(0)$. Take $G(1)=\{r,s\}$, $G(0)=\{a,b,c\}$, with restrictions $r\mapsto a$, $s\mapsto b$. Let $F$ be the subpresheaf with $F(1)=\{r\}$ and $F(0)=\{a,b\}$, and let $F'=e$ be the terminal presheaf. The unique map $F\to e$ gives the pushout
\[
 G'(1)=\{*,s\},\qquad G'(0)=\{*,c\},
 \qquad *\mapsto*,\quad s\mapsto*.
\]

The category $X/G$ has five objects and two nonidentity arrows,
\[
 (0,a)\longrightarrow(1,r),\qquad
 (0,b)\longrightarrow(1,s),
\]
with $(0,c)$ isolated. Its open subcategory $X/F$ has three objects: the first arrow and the isolated object $(0,b)$. The category $X/e$ is $[1]$. In the pushout, $(1,r)$ is identified with $(1,*)$, and both $(0,a)$ and $(0,b)$ with $(0,*)$. The resulting category has exactly the two arrows
\[
 (1,*)\longleftarrow(0,*)\longrightarrow(1,s)
\]
and the isolated object $(0,c)$. This is $X/G'$. There are no further nonidentity compositions to add.

The two complementary closed subcategories are both the discrete category on $(1,s)$ and $(0,c)$. The new attachment changes the crossing arrow into $(1,s)$ but leaves the complement itself unchanged, exactly as Proposition 15 predicts. Their nerves give a homotopy pushout. On connected components the four corners have respectively $2,3,1,2$ components, and the last set is the pushout of the first three component sets. This also checks concretely that the word ``cartesian'' for the presheaf square does not mean that either vertical map is an equivalence: the left vertical map merges two components, and the right one merges the corresponding two components after attachment.
\end{dcomment}


\textbf{Corollaire :} (\footnote[111]{\textbf{NB} Si on veut seulement $i \in W_X \Longrightarrow i' \in W_X$, $p \in W_X \Longrightarrow q \in W_X$, il suffit que l’on ait W(7 bis) \textit{faible}, i.e. $i_0 \in W \Longrightarrow i \in W$ dans l’énoncé de W(7 bis), p. 76 (Loc(6$'$a) avec les notations de XVI).}) \textit{On a} (moyennant W(7 bis)) \textit{les équivalences}
\[
\begin{aligned}
i \in W_X \quad &\Longleftrightarrow \quad i' \in W_X \\
p \in W_X \quad &\Longleftrightarrow \quad q \in W_X,
\end{aligned}
\]
\textit{où}
\[
W_X \subseteq \mathit{Fl}(X^\wedge)
\]
\textit{est défini par (cf. Pursuing Stacks \ldots)}
\[
u \in W_X \qquad\qquad \Longleftrightarrow \qquad\qquad X/u : X/F \longrightarrow X/G \quad \textit{est dans } W.
\]

Ainsi, dans une catégorie $X^\wedge$, les sommes amalgamées impliquant un monomorphisme se comportent de façon parfaite du point de vue homotopique, plus exactement pour les localiseurs $W_X$, pourvu \msp{page 169} seulement que $W$ satisfasse à W(7 bis). Notamment, les monomorphismes dans $X^\wedge$ sont des ‘$W_X$-cofibrations’, en un sens similaire de celui de la définition 2 c) (pages 99, 100 \texttt{[plutôt 109, 110]}).

\texttt{[page 170 n’existe pas.]}


\dgrblock{dgrb481}{From attachment stability to model structures}
\begin{dcomment}{status}{p169}{From attachment stability to model structures}
The final assertion concerns a class of cofibrations in the sense developed in this chapter: pushing a $W_X$-equivalence out along a monomorphism again gives a $W_X$-equivalence. This preservation property survives for ordinary homotopy types. The two reverse implications displayed in the corollary are different assertions and depend on the strong W(7 bis) hypothesis. The earlier failure of pushout reflection applies here as well after representing spaces by simplicial presheaves.

It would go beyond the text to infer that, for every small $X$, the class $W_X$ is the weak-equivalence class of a model structure whose cofibrations are all monomorphisms. Already $X=e$ shows the missing condition. For $W=W_\infty$, a presheaf is a set and $X/F$ its discrete category, so $W_X$ consists of bijections. If a model structure on sets had these weak equivalences and all monomorphisms as cofibrations, every trivial fibration would be a bijection. Factoring a noninjective map as a cofibration followed by a trivial fibration would make it injective, a contradiction. Attachment stability alone does not supply the factorization axiom.

The later theory identifies the additional hypotheses. For an accessible fundamental localizer and a $W$-local test category, the induced presheaf equivalences admit a cofibrantly generated model structure with monomorphisms as cofibrations; this is \citet[Corollary 4.2.18]{Cisinski2006}. In the ordinary case, the local test condition is characterized by the requirement that maps with the right lifting property against all monomorphisms be weak equivalences; see \citet[Theorem 4.1.19]{Cisinski2006}. The general existence theorem for accessible presheaf localizers, Theorem 1.4.3 of the same work, supplies the factorizations. The present calculation establishes the attachment mechanism on which that construction depends, while the test-category condition ensures that it models the intended homotopy theory.
\end{dcomment}
\begin{otherlanguage}{english}
\dgroverview{XIIb-b}{Chapter XII \textperiodcentered\ \textup{XIIb-b} \textperiodcentered\ Overview}
\selectlanguage{english}
\noindent\textbf{Chapter XII, section 5, manuscript pages 115--169.}
\medskip

These pages develop a concrete calculus for deciding when an ordinary gluing diagram of categories represents a homotopy gluing. Its central object is a comparison map. Given a strictly commuting square with initial corner $X_0$, adjacent corners $X_1,X_2$, and terminal corner $Y$, Grothendieck constructs a category $\mathcal X_Q$ from the three-category span and a canonical functor
\[
                    \varphi_Q:\mathcal X_Q\longrightarrow Y.
\]
The square is $W$-homotopically cocartesian when this particular functor belongs to the chosen weak-equivalence class $W$. The definition remembers the cocone, not only its four vertex homotopy types. The work that follows explains why the construction is independent of its two orientations, how smoothness can test the comparison on fibres, and why monomorphisms of presheaves provide a reliable supply of such squares.

The opening argument concerns a different but related exactness problem. Total $W$-asphericity here is the requirement that, for every $p:X\to Y$, the category of elements of the ordinary left Kan extension $p_!1_X$ be $W$-aspherical. It is the definition of this chapter, explicitly distinguished by the editors from the earlier use of the same phrase in \emph{Pursuing Stacks}. $W_0$-properness, or $W$-properness under W(5), allows Grothendieck to compute
\[
 p_!1_X\cong\operatorname*{colim}_{i\in I}p_{i!}1_{X_i}
\]
for a family $X\to I$. A terminal object of $I$ reduces this colimit to the last fibre. The commentary supplies the generators-and-relations proof on components and the precise functor on comma categories requested in footnote~83. It also proves that total $W_\infty$-asphericity is not invariant under weak homotopy equivalence. The properness hypothesis controls ordinary Kan extensions; equivalence of classifying spaces alone does not do that work.

The smooth integral of $X_1\leftarrow X_0\to X_2$ keeps the three sets of objects disjoint. Its mixed arrows run from the outer fibres to the middle fibre, with
\[
 \operatorname{Hom}_{\mathcal X_Q}(a,x)
       =\operatorname{Hom}_{X_i}(a,\alpha_i x)
       \qquad(a\in X_i,\ x\in X_0).
\]
Its universal property therefore includes natural transformations $f_i\alpha_i\Rightarrow f_0$. They need not be invertible. Strict commutativity of the given square supplies identity transformations and hence the canonical $\varphi_Q$. Dividing the integral into two full closed pieces gives a genuine categorical pushout and pullback. Each piece retracts onto its outer fibre by an adjunction. The nerve of this closed cover proves directly that the total category represents the homotopy pushout of classifying spaces. The general theorem is Thomason's; the bibliography distinguishes verification of the original metadata from verification of an independent primary proof \citep{Thomason1979,Cegarra2020}.

Even the smallest examples explain why this construction is needed. For the span $e\leftarrow\{a,b\}\to e$, the integral is a four-object category whose nerve is a circle. Its ordinary categorical pushout is a point. The ordinary pushout of the same diagram in the homotopy category is also a point: it retains equality of components but loses the two paths implementing compatibility. The integral retains those paths as transition arrows. This is the concrete content, on these pages, of the insistence on forming diagram categories before localization.

The excision hypotheses require particular care. Proposition~13 and its cube corollary are valid conditional assertions under the full W(7 bis) and W(7) printed in the manuscript. Those axioms include reflection statements stronger than the usual preservation of weak equivalences by homotopy pushout. They fail for $W_\infty$. The commentary proves this with a connected acyclic finite complex having nontrivial fundamental group. Its suspension is contractible, so the union of its two cones and both cones are contractible although their intersection is not. A map of this closed-cover square to the point square violates the converse cube implication. The usual forward gluing theorem survives. A normalized-chain calculation also explains exactly why the complete three-way implication works for integral homology equivalences. The earlier chapter's claim about full W(7) for $W_\infty$ is therefore not a valid justification of the converses in the present portion.

The proper integral reverses the transition arrows. Grothendieck compares the two constructions by the functorial resolution
\[
                         C\longleftarrow\Sigma(C)\longrightarrow C^{\mathrm{op}}.
\]
Both projections are aspherical, and the same comparison can be integrated. Proposition~14 proves agreement of the four versions obtained from smooth or proper integration and from the span or its opposite. The commentary writes out the chosen twisted-arrow convention and completes the deferred transitivity verification. The natural comparison with the opposite operation is involutive; its naturality makes the square of four integration comparisons commute. The indexing category must keep its two paths distinct until they reach $Y$. Replacing that free diamond by a commutative square would erase precisely the information the proof is using.

The question on manuscript pages 145--146 asks whether a bad ordinary pushout can become a good homotopy cocone after mapping to another target. We answer it with finite posets. A collapsible disk with seven triangles maps to the tetrahedral two-sphere so that all associated comma categories are nonempty and connected. Its face poset $E$ embeds as a sieve in a contractible mapping-cylinder poset $B$. The categorical pushout
\[
                         P=B\amalg_E e
\]
has two incomparable minimal vertices below the face poset of that sphere. Its nerve is their join with the subdivided sphere, hence has classifying space $S^3$. The span $B\leftarrow E\to e$ has contractible integral. Consequently the comparison to $P$ is not a weak equivalence, whereas its composite with $P\to e$ is. The proof includes the categorical pushout universal property; a separate exact finite calculation checks the collapses, comma connectivity, and rational homology. This is a worked development, with no claim of priority. The same disk-to-sphere map supplies the promised counterexample to homotopy invariance of total asphericity.

Theorem~4 turns from construction to recognition. Its forward assertion says that smoothness of the three maps $f_i$ implies smoothness of $\varphi_Q$. The reverse assertion for the middle map retains strong W(7 bis). Two slice calculations identify the strict fibres relevant to local smoothness with integrals of lift categories. Footnote~107 correctly observes that the apparent square in the original proof is not strictly commutative. The repair uses the split fibration of the integral over the fork and its aspherical fibres. This proves the forward direction without the strong excision axiom. An explicit contraction removes the same unnecessary hypothesis from the open-cover corollaries. The final unconditional sentence after Corollary~6, however, omits essential outer smoothness assumptions; a four-object lift category with an isolated component gives a finite counterexample.

The last part returns to ordinary category theory. A full sieve $U\subset X$ and its complement $T$ are glued by the presheaves $\operatorname{Hom}_X(-,t)$ on $U$. Proposition~15 says that a square along $U\subset X$ is a pushout exactly when it identifies the complementary categories and transports these crossing-arrow presheaves by left Kan extension. The connected factorization categories in its statement are the fibres of that comparison of presheaves. This gives a full proof of the criterion and explains why an agreement on objects and strata would be insufficient.

For a presheaf $F$ on $X$, the formula
\[
 N(X/F)_n=\coprod_{x_0\to\cdots\to x_n}F(x_n)
\]
then makes the final gluing theorem transparent. It proves preservation of colimits by the category-of-elements construction; the relative construction into $\Cat/X$ also preserves all small limits, the known strengthening recorded by \citet[Proposition~3.2.12 and Corollary~3.2.13]{Cisinski2006}. A pushout of presheaves along a monomorphism becomes a simplicial pushout along a monomorphism. Thus the positive homotopy-pushout clause of Theorem~5 holds for every fundamental localizer. Its subsequent reflection implications still require the strong hypothesis. Model structures require further conditions, supplied later by accessibility and local test categories; attachment stability alone is insufficient.

The later derivator definition identifies $\varphi_Q$ with the terminal component of the coherent left-Kan-extension counit \citep[Definition~3.9]{Groth2013}. The comparison with $\infty$-categories requires another localization: a total category models the categorical colimit after its transport arrows are inverted, while classifying spaces invert all arrows \citep[Corollary~3.3.4.3]{Lurie2009HTT}. These formulations clarify the target of the present calculations without replacing their elementary categorical proofs. The equivalence of the two integration conventions and the strict-cocone question are resolved, as are the specified coherence and local-smoothness gaps. The opening request for a sufficiently large class of cofibrations closed under the stated interval operation is not settled here in its full generality. We record that limitation without declaring it a present-day open problem.
\end{otherlanguage}
\dgrchunkend
